\section{AMC 10 10B Fall 2021}
        \begin{enumerate}
        	\item (\textit{AMC 10 10B Fall 2021}) What is the value of $1234+2341+3412+4123?$
$(\textbf{A})\: 10{,}000\qquad(\textbf{B}) \: 10{,}010\qquad(\textbf{C}) \: 10{,}110\qquad(\textbf{D}) \: 11{,}000\qquad(\textbf{E}) \: 11{,}110$

 \textbf{Answer: }E



\textbf{Solution 1}

We see that $1, 2, 3,$ and $4$ each appear in the ones, tens, hundreds, and thousands digit exactly once. Since $1+2+3+4=10$, we find that the sum is equal to 
 \[10\cdot(1+10+100+1000)=\boxed{\textbf{(E)} \: 11{,}110}.\]Note that it is equally valid to manually add all four numbers together to get the answer.

 



\textbf{Solution 2}

We have 
 \[1234 + 2341 + 3412 + 4123 = 1111 \left( 1 + 2 + 3 + 4 \right) = \boxed{\textbf{(E)} \: 11{,}110}.\]~Steven Chen (www.professorchenedu.com)

 



\textbf{Solution 3}

We see that the units digit must be $0$, since $4+3+2+1$ is $0$. But every digit from there, will be a $1$ since we have that each time afterwards, we must carry the $1$ from the previous sum. The answer choice that satisfies these conditions is $\boxed{\textbf{(E)} \: 11{,}110}$.

 



\textbf{Solution 4 (Brute Force)}

We can simply add the numbers. $1234 + 2341 + 3412 + 4123 = 11110 \implies \boxed{\textbf{(E)}}$.

 Note: Although this would not take terribly long, it is not recommended to do this in a real contest.

 



\textbf{Solution 5 (Primary School But Also Similar to Solution 4)}

We can simply just add each place value. 

 $1000+2000+3000+4000=10,000$
$+200+300+400+100=1,000$
$+30+40+10+20=100$
$+4+1+2+3=10$
$10,000+1,000+100+10=\boxed{\textbf{(E)} \: 11{,}110}$.

 


	\item (\textit{AMC 10 10B Fall 2021}) What is the area of the shaded figure shown below?

                    \begin{figure}[H]
                        \centering
                        \includegraphics[width=0.5\textwidth]{Problem/AMC_10_10B_Fall_2021_2_Problem_Figure_1.png}
                    \end{figure}
                    
$\textbf{(A)}\: 4\qquad\textbf{(B)} \: 6\qquad\textbf{(C)} \: 8\qquad\textbf{(D)} \: 10\qquad\textbf{(E)} \: 12$

 \textbf{Answer: }B



\textbf{Solution 1 (Area Addition)}

The line of symmetry divides the shaded figure into two congruent triangles, each with base $3$ and height $2.$

 Therefore, the area of the shaded figure is 
 \[2\cdot\left(\frac12\cdot3\cdot2\right)=2\cdot3=\boxed{\textbf{(B)} \: 6}.\]~MRENTHUSIASM ~Wilhelm Z

 



\textbf{Solution 2 (Area Subtraction)}

To find the area of the shaded figure, we subtract the area of the smaller triangle (base $4$ and height $2$) from the area of the larger triangle (base $4$ and height $5$): 
 \[\frac12\cdot4\cdot5-\frac12\cdot4\cdot2=10-4=\boxed{\textbf{(B)} \: 6}.\]~MRENTHUSIASM ~Steven Chen (www.professorchenedu.com) 

 



\textbf{Solution 3 (Shoelace Theorem)}

The consecutive vertices of the shaded figure are $(1,0),(3,2),(5,0),$ and $(3,5).$ By the \href{https://artofproblemsolving.com/wiki/index.php?title=Shoelace\_Theorem}{Shoelace Theorem}, the area is 
 \[\frac12\cdot|(1\cdot2+3\cdot0+5\cdot5+3\cdot0)-(0\cdot3+2\cdot5+0\cdot3+5\cdot1)|=\frac12\cdot12=\boxed{\textbf{(B)} \: 6}.\]~Taco12 ~I-AM-DA-KING

 



\textbf{Solution 4 (Pick's Theorem)}

We have $4$ lattice points in the interior and $6$ lattice points on the boundary. By \href{https://artofproblemsolving.com/wiki/index.php?title=Pick\%27s\_Theorem}{Pick's Theorem}, the area of the shaded figure is 
 \[4+\frac{6}{2}-1 = 4+3-1 = \boxed{\textbf{(B)} \: 6}.\]~danprathab

 


	\item (\textit{AMC 10 10B Fall 2021}) The expression $\frac{2021}{2020} - \frac{2020}{2021}$ is equal to the fraction $\frac{p}{q}$ in which $p$ and $q$ are positive integers whose greatest common divisor is ${ }1$. What is $p?$
$\textbf{(A)}\: 1\qquad\textbf{(B)} \: 9\qquad\textbf{(C)} \: 2020\qquad\textbf{(D)} \: 2021\qquad\textbf{(E)} \: 4041$

 \textbf{Answer: }E



\textbf{Solution 1}

We write the given expression as a single fraction: 
 \[\frac{2021}{2020} - \frac{2020}{2021} = \frac{2021\cdot2021-2020\cdot2020}{2020\cdot2021}\] by cross multiplication. Then by factoring the numerator, we get 
 \[\frac{2021\cdot2021-2020\cdot2020}{2020\cdot2021}=\frac{(2021-2020)(2021+2020)}{2020\cdot2021}.\] The question is asking for the numerator, so our answer is $2021+2020=4041,$ giving $\boxed{\textbf{(E) }4041}$.

 



\textbf{Solution 2}

Denote $a = 2020$. Hence,
 \begin{align*} \frac{2021}{2020} - \frac{2020}{2021} & = \frac{a + 1}{a} - \frac{a}{a + 1} \\ & = \frac{\left( a + 1 \right)^2 - a^2}{a \left( a + 1 \right)} \\ & = \frac{2 a + 1}{a \left( a + 1 \right)} . \end{align*}


 We observe that ${\textrm gcd} \left( 2a + 1 , a \right) = 1$ and ${\textrm gcd} \left( 2a + 1 , a + 1 \right) = 1$.

 Hence, ${\textrm gcd} \left( 2a + 1 , a \left( a + 1 \right) \right) = 1$.

 Therefore, $p = 2 a + 1 = 4041$.

 Therefore, the answer is $\boxed{\textbf{(E) }4041}$.

 



\textbf{Solution 3}

Turning term 1 to a fraction:$\frac{2021}{2020}-\frac{2020}{2021}=1+\frac{1}{2020}-\frac{2020}{2021}$

 Subtracting the last term from the first term gives us:$\frac{1}{2021}+\frac{1}{2020}$

 Doing some simple cross multiplication, you get $\frac{4041}{2020\cdot2021}$, here you can see the numerator is $\boxed{\textbf{(E)}4041}$.

 


	\item (\textit{AMC 10 10B Fall 2021}) At noon on a certain day, Minneapolis is $N$ degrees warmer than St. Louis. At $4{:}00$ the temperature in Minneapolis has fallen by $5$ degrees while the temperature in St. Louis has risen by $3$ degrees, at which time the temperatures in the two cities differ by $2$ degrees. What is the product of all possible values of $N?$
$\textbf{(A)}\: 10\qquad\textbf{(B)} \: 30\qquad\textbf{(C)} \: 60\qquad\textbf{(D)} \: 100\qquad\textbf{(E)} \: 120$

 \textbf{Answer: }C



\textbf{Solution 1 (Two Variables)}

At noon on a certain day, let $M$ and $L$ be the temperatures (in degrees) in Minneapolis and St. Louis, respectively. It follows that $M=L+N.$

 At $4{:}00,$ we get
 \begin{align*} |(M-5)-(L+3)| &= 2 \\ |M-L-8| &= 2 \\ |N-8| &= 2. \end{align*}
We have two cases:

 
            \begin{enumerate} 
        \item If $N-8=2,$ then $N=10.$

        \item If $N-8=-2,$ then $N=6.$

            \end{enumerate}
            Together, the product of all possible values of $N$ is $10\cdot6=\boxed{\textbf{(C)} \: 60}.$




\textbf{Solution 2 (One Variable)}

At noon on a certain day, the difference of temperatures in Minneapolis and St. Louis is $N$ degrees.

 At $4{:}00,$ the difference of temperatures in Minneapolis and St. Louis is $N-8$ degrees.

 It follows that 
 \[|N-8|=2.\]We continue with the casework in Solution 1 to get the answer $\boxed{\textbf{(C)} \: 60}.$




\textbf{Solution 3}

Let the temperature in Minneapolis be $N$. Then St. Louis is $0$ degrees. $N - 5 = 3 \pm 2$, $N = -10, -6$, $-10 \cdot -6 = \boxed{\textbf{(C)} \: 60}$.

 


	\item (\textit{AMC 10 10B Fall 2021}) Let $n=8^{2022}$. Which of the following is equal to $\frac{n}{4}?$
$\textbf{(A)}\: 4^{1010}\qquad\textbf{(B)} \: 2^{2022}\qquad\textbf{(C)} \: 8^{2018}\qquad\textbf{(D)} \: 4^{3031}\qquad\textbf{(E)} \: 4^{3032}$

 \textbf{Answer: }E



\textbf{Solution 1}

We have 
 \[n=8^{2022}=  \left(8^\frac{2}{3}\right)^{3033}=4^{3033}.\]Therefore, 
 \[\frac{n}4=\boxed{\textbf{(E)} \: 4^{3032}}.\]~kingofpineapplz

 



\textbf{Solution 2}

The requested value is 
 \[\frac{8^{2022}}{4} = \frac{2^{6066}}{4} = \frac{2^{6066}}{2^2} = 2^{6064} = \boxed{\textbf{(E)} \: 4^{3032}}.\]~NH14

 



\textbf{Solution 3}

If we rewrite everything in powers of 2, we get:$n = 8^{2022} = (2^{3})^{2022} = 2^{6064} = (4^{\frac{1}{2}})^{6064} = 4^{\frac{6064}{2}} = \boxed{\textbf{(E)} \: 4^{3032}}.$

 - abed\_nadir (youtube.com/@indianmathguy)

 



\textbf{Solution 4}

$n = 8^{2022}$
${8^{2022}}/{4} = 2 * 8^{2021} = 2 * 2^{6063} = 2^{6064} = \boxed{\textbf{(E)} \: 4^{3032}}.$

 -sphericalfox

 


	\item (\textit{AMC 10 10B Fall 2021}) The least positive integer with exactly $2021$ distinct positive divisors can be written in the form $m \cdot 6^k$, where $m$ and $k$ are integers and $6$ is not a divisor of $m$. What is $m+k?$
$(\textbf{A})\: 47\qquad(\textbf{B}) \: 58\qquad(\textbf{C}) \: 59\qquad(\textbf{D}) \: 88\qquad(\textbf{E}) \: 90$

 \textbf{Answer: }B



\textbf{Solution 1}

Let this positive integer be written as $p_1^{e_1}\cdot p_2^{e_2}$. The number of factors of this number is therefore $(e_1+1) \cdot (e_2+1)$, and this must equal 2021. The prime factorization of 2021 is $43 \cdot 47$, so $e_1+1 = 43 \implies e_1=42$ and $e_2+1=47\implies e_2=46$. To minimize this integer, we set $p_1 = 3$ and $p_2 = 2$. Then this integer is $3^{42} \cdot 2^{46} = 2^4 \cdot 2^{42} \cdot 3^{42} = 16 \cdot 6^{42}$.Now $m=16$ and $k=42$ so $m+k = 16 + 42 = \boxed{\textbf{(B) }58}$




\textbf{Solution 2}

Recall that $6^k$ can be written as $2^k \cdot 3^k$. Since we want the integer to have $2021$ divisors, we must have it in the form $p_1^{42} \cdot p_2^{46}$, where $p_1$ and $p_2$ are prime numbers. Therefore, we want $p_1$ to be $3$ and $p_2$ to be $2$. To make up the remaining $2^4$, we multiply $2^{42} \cdot 3^{42}$ by $m$, which is $2^4$ which is $16$. Therefore, we have $42 + 16 = \boxed{\textbf{(B) }58}$




\textbf{Solution 3}

If a number has prime factorization $p_1^{k_1} p_2^{k_2} \cdots p_m^{k_m}$, then the number of distinct positive divisors of this number is $\left( k_1 + 1 \right) \left( k_2 + 1 \right) \cdots \left( k_m + 1 \right)$.

 We have $2021 = 43 \cdot 47$.Hence, if a number $N$ has 2021 distinct positive divisors, then $N$ takes one of the following forms: $p_1^{2020}$, $p_1^{42} p_2^{46}$.

 Therefore, the smallest $N$ is $3^{42} 2^{46} = 2^4 \cdot 6^{42} = 16 \cdot 6^{42}$.

 Therefore, the answer is $\boxed{\textbf{(B) }58}$.

 


	\item (\textit{AMC 10 10B Fall 2021}) Call a fraction $\frac{a}{b}$, not necessarily in the simplest form, \textit{special}, if $a$ and $b$ are positive integers whose sum is $15$. How many distinct integers can be written as the sum of two, not necessarily different, special fractions?


 $\textbf{(A)}\ 9 \qquad\textbf{(B)}\  10 \qquad\textbf{(C)}\  11 \qquad\textbf{(D)}\ 12 \qquad\textbf{(E)}\ 13$

 \textbf{Answer: }C



\textbf{Solution 1}

The special fractions are 
 \[\frac{1}{14},\frac{2}{13},\frac{3}{12},\frac{4}{11},\frac{5}{10},\frac{6}{9},\frac{7}{8},\frac{8}{7},\frac{9}{6},\frac{10}{5},\frac{11}{4},\frac{12}{3},\frac{13}{2},\frac{14}{1}.\]We rewrite them in the simplest form: 
 \[\frac{1}{14},\frac{2}{13},\frac{1}{4},\frac{4}{11},\frac{1}{2},\frac{2}{3},\frac{7}{8},1\frac{1}{7},1\frac{1}{2},2,2\frac{3}{4},4,6\frac{1}{2},14.\]Note that two unlike fractions in the simplest form cannot sum to an integer. So, we only consider the fractions whose denominators appear more than once: 
 \[\frac{1}{4},\frac{1}{2},1\frac{1}{2},2,2\frac{3}{4},4,6\frac{1}{2},14.\]For the set $\{2,4,14\},$ two elements (not necessarily different) can sum to $4,6,8,16,18,28.$

 For the set $\left\{\frac{1}{2},1\frac{1}{2},6\frac{1}{2}\right\},$ two elements (not necessarily different) can sum to $1,2,3,7,8,13.$

 For the set $\left\{\frac{1}{4},2\frac{3}{4}\right\},$ two elements (not necessarily different) can sum to $3.$

 Together, there are $\boxed{\textbf{(C)}\  11}$ distinct integers that can be written as the sum of two, not necessarily different, special fractions: 
 \[1,2,3,4,6,7,8,13,16,18,28.\]~KingRavi ~samrocksnature ~Wilhelm Z ~MRENTHUSIASM ~Steven Chen (www.professorchenedu.com)

 



\textbf{Solution 2}

Let $a=15-b,$ so the special fraction is 
 \[\frac ab = \frac{15-b}{b} = \frac{15}{b}-1.\]We can ignore the $-1$ part and only focus on $\frac{15}{b}.$

 The integers are $\frac{15}{1},\frac{15}{3},\frac{15}{5},$ which are $15,5,3,$ respectively. We get $30,20,18,10,8,6$ from this group of numbers.

 The halves are $\frac{15}{2},\frac{15}{6},\frac{15}{10},$ which are $7\frac12,2\frac12,1\frac12,$ respectively. We get $15,10,9,5,4,3$ from this group of numbers.

 The quarters are $\frac{15}{4},\frac{15}{12},$ which are $3\frac34,1\frac14,$ respectively. We get $5$ from this group of numbers.

 Note that $10$ and $5$ each appear twice. Therefore, the answer is $\boxed{\textbf{(C)}\  11}.$




\textbf{Solution 3 (Casework)}

We split this up into two cases:

 \textbf{Case 1: integer + integer}

 The whole numbers we have are $\frac{10}{5}$ (or $2$), $\frac{12}{3}$ (or $4$), and $\frac{14}{1}$ (or $14$). There are $\dbinom{3}{2}=3$ ways to choose different-numbered pairs and $3$ ways to choose the same-numbered pairs. So, $3+3=6$.

 \textbf{Case 2: fraction + fraction}

 The fractions we have are $\frac{5}{10}$ (or $\frac{1}{2}$), $\frac{9}{6}$ (or $\frac{3}{2}$), and $\frac{13}{2}$. Similarly, there are $\dbinom{3}{2}=3$ ways to choose different-numbered pairs and $3$ ways to choose the same-numbered pairs. So, $3+3=6$.

 Thus, $6+6=12$.

 So now you would just go ahead and innocently choose $\textbf{(D) }12$, right? No! We overcounted $8$, as $\frac{13}{2}+\frac96=\frac{12}{3}+\frac{12}{3}=8$. Therefore, the correct answer is actually $12-1=\boxed{\textbf{(C)}\  11}$.

 


	\item (\textit{AMC 10 10B Fall 2021}) The greatest prime number that is a divisor of $16{,}384$ is $2$ because $16{,}384 = 2^{14}$. What is the sum of the digits of the greatest prime number that is a divisor of $16{,}383$?


 $\textbf{(A)} \: 3\qquad\textbf{(B)} \: 7\qquad\textbf{(C)} \: 10\qquad\textbf{(D)} \: 16\qquad\textbf{(E)} \: 22$

 \textbf{Answer: }C



\textbf{Solution}

We have
 \begin{align*} 16383 & = 2^{14} - 1 \\ & = \left( 2^7 + 1 \right) \left( 2^7 - 1 \right) \\ & = 129 \cdot 127 \\ \end{align*}


 Since $129$ is composite, $127$ is the largest prime which can divide $16383$. The sum of $127$'s digits is $1+2+7=\boxed{\textbf{(C) }10}$.

 \textbf{Note} Note that you can quickly tell that $2^7 -1$ is prime because it is a \href{https://artofproblemsolving.com/wiki/index.php?title=Mersenne\_prime}{Mersenne prime}.

 


	\item (\textit{AMC 10 10B Fall 2021}) The knights in a certain kingdom come in two colors. $\frac{2}{7}$ of them are red, and the rest are blue. Furthermore, $\frac{1}{6}$ of the knights are magical, and the fraction of red knights who are magical is $2$ times the fraction of blue knights who are magical. What fraction of red knights are magical?


 $\textbf{(A) }\frac{2}{9}\qquad\textbf{(B) }\frac{3}{13}\qquad\textbf{(C) }\frac{7}{27}\qquad\textbf{(D) }\frac{2}{7}\qquad\textbf{(E) }\frac{1}{3}$

 \textbf{Answer: }C



\textbf{Solution 1}

Let $k$ be the number of knights: then the number of red knights is $\frac{2}{7}k$ and the number of blue knights is $\frac{5}{7}k$.

 Let $b$ be the fraction of blue knights that are magical - then $2b$ is the fraction of red knights that are magical. Thus we can write the equation $b \cdot \frac{5}{7}k + 2b \cdot \frac{2}{7}k = \frac{k}{6}\implies \frac{5}{7}b + \frac{4}{7}b = \frac{1}{6}$$\implies \frac{9}{7}b = \frac{1}{6} \implies b=\frac{7}{54}$

 We want to find the fraction of red knights that are magical, which is $2b = \frac{7}{27} = \boxed{\textbf{(C) }\frac{7}{27}}$




\textbf{Solution 2}

We denote by $p$ the fraction of red knights who are magical.

 Hence,
 \[ \frac{1}{6} = \frac{2}{7} p + \left( 1 - \frac{2}{7} \right) \frac{p}{2} . \]

 By solving this equation, we get $p = \frac{7}{27}$.

 Therefore, the answer is $\boxed{\textbf{(C) }\frac{7}{27}}$.

 



\textbf{Solution 3}

Suppose there are $42$ knights. Then there are $\frac{2}{7}\cdot 42=12$ red knights and $42-12=30$ blue knights. 

 There are $\frac{1}{6}\cdot 42=7$ magical knights. Let $x$ be the number of red magical knights. 

 Then we have the equation $\frac{x}{12}=2\cdot \frac{7-x}{30}=\frac{7-x}{15}$. Solving, we get $x=84/27$. 

 Plugging $x$ back gives us the fraction of red magical knights, which is $\frac{\frac{84}{27}}{12}=\frac{84}{27}\cdot \frac{1}{12}=\boxed{\textbf{(C) }\frac{7}{27}}$.

 


 



\textbf{Solution 4}

Similar to Solution 3, let the total number of knights be $42$ (We can say this because the fractions we know have denominators of $7$ and $6$ which have a LCM of $42$)

 This means that there are $\dfrac{2}{7} \times 42 = 12$ red knights meaning there are $30$ blue knights

 This also means that there are $\dfrac{1}{6} \times 42 = 7$ magical knights. This is important information for our next step.

 We are given that the fraction of magical red knights is twice the fraction of magical blue knights. As this problem is dealing with the fraction being twice are large and is asking for the fraction, let's just deal with the variable being a fraction. Let $\dfrac{p}{q}$ be the fraction of red magical knights. This means that $\dfrac{p}{2q}$ is the fraction of blue magical knights.

 We know that total number of magical knights from what we figured out before is $7$, so let's use the fraction we figured and make the equation:

 $(12 \times \dfrac{p}{q}) + (30 \times \dfrac{p}{2q})= 7$

 Solving this gives us $\dfrac{p}{q}=\dfrac{7}{27}$ or the fraction of red magical knights.

 Hence, the answer is $\boxed{\dfrac{7}{27}}$




\textbf{Solution 5}

This solution is similar to Solution $2$.

 Suppose $x$ is the fraction of red knights that are magical. Then $\frac{x}{2}$ is the fraction of blue knights that are magical. We can then create an equation in terms of $x$.

 
 \[ \frac{2}{7}x+\frac{5}{7}(\frac{x}{2})=\frac{1}{6} \]

 Solving, we get that $x=\frac{7}{27}$.

 The answer is $\boxed{\textbf{(C) }\frac{7}{27}}$.

 


	\item (\textit{AMC 10 10B Fall 2021}) Forty slips of paper numbered $1$ to $40$ are placed in a hat. Alice and Bob each draw one number from the hat without replacement, keeping their numbers hidden from each other. Alice says, "I can't tell who has the larger number." Then Bob says, "I know who has the larger number." Alice says, "You do? Is your number prime?" Bob replies, "Yes." Alice says, "In that case, if I multiply your number by $100$ and add my number, the result is a perfect square. " What is the sum of the two numbers drawn from the hat?


 $\textbf{(A) }27\qquad\textbf{(B) }37\qquad\textbf{(C) }47\qquad\textbf{(D) }57\qquad\textbf{(E) }67$

 \textbf{Answer: }A



\textbf{Solution 1}

Let Alice have the number A, Bob B.When Alice says that she can't tell who has the larger number, it means that $A$ cannot equal $1$. Therefore, it makes sense that Bob has $2$ because he now knows that Alice has the larger number. $2$ is also prime. The last statement means that $200+A$ is a perfect square. The three squares in the range $200-300$ are $225$, $256$, and $289$. So, $A$ could equal $25$, $56$, or $89$, so $A+B$ is $27$, $58$, or $91$, of only $\boxed{\textbf{(A) }27}$ is below $40$, therefore it is the correct answer since the maximum number one can have is 40.

 --stjwyl

 Minor edit ~SIGMAMATHEMATICIAN

 



\textbf{Solution 2}

Denote by $A$ and $B$ the numbers drawn by Alice and Bob, respectively.

 Alice's sentence “I can't tell who has the larger number.” implies $A \in \left\{ 2 , \cdots , 39 \right\}$.

 Bob's sentence “I know who has the larger number.” implies $B \in \left\{ 1 , 2 , 39, 40 \right\}$.

 Their subsequent conversation that $B$ is prime implies $B = 2$.

 Then, Alice's next sentence  “In that case, if I multiply your number by 100 and add my number, the result is a perfect square.” implies $200 + A$ is a perfect square.Hence, $A = 25$.

 Therefore, the answer is $\boxed{\textbf{(A) }27}$.

 



\textbf{Solution 3 - Guessing those Squares}

We see that $225$ is one such square that works. Bob gets $2$ and Alice gets $25$ which is valid. Thus, $2 + 25 = 27.$ So $\boxed{\textbf{(A) }27}$ is our answer.

 -D1r

 



\textbf{Solution 4 (Ambiguity, Sanity Checks)}

The MAA knows this problem is ambiguous, as the interpretation of what might be viewed as the "middle" of the set could be different for all. Then, they would have specifically designed the answer choices such that one choice is guaranteed to be correct while all others are unobtainable. 

 So, for a sanity check and our general satisfaction, we will use the set $\mathcal{A} \in \{5,6,7,\dots,34,35 \}$ as all the possible values of $\mathcal{A}$ and $\mathcal{B} \in \{1,2,3,4,36,37,38,39,40 \}$. Alice notes that Bob's number is prime. Therefore, $\mathcal{B} \in \{2,3,4,37\}$. Now, say Alice's number is $A$. Then, each of $200+A$, $300+A$, and $3700+A$ is prime. Note that in each case, $A = 25, 24, 21$ respectively. Then, $A+B = 27, 27, 58$ respectively. Only $\boxed{27}$ is in the answer choices, so that would be our answer.

 \textbf{Sidenote} Note that Bob's statement (that he knows who won/lost) comes after Alice tells him that she doesn't know who won. Since Alice doesn't know who won, Bob knows that she didn't get draw a $1$ (or a $40$), which tells him that his $2$ must be lower than Alice's number. 

 


	\item (\textit{AMC 10 10B Fall 2021}) A regular hexagon of side length $1{ }$ is inscribed in a circle. Each minor arc of the circle determined by a side of the hexagon is reflected over that side. What is the area of the region bounded by these $6$ reflected arcs?


 $(\textbf{A})\: \frac{5\sqrt{3}}{2} - \pi\qquad(\textbf{B}) \: 3\sqrt{3}-\pi\qquad(\textbf{C}) \: 4\sqrt{3}-\frac{3\pi}{2}\qquad(\textbf{D}) \: \pi - \frac{\sqrt{3}}{2}\qquad(\textbf{E}) \: \frac{\pi + \sqrt{3}}{2}$

 \textbf{Answer: }B



\textbf{Solution 1}


                    \begin{figure}[H]
                        \centering
                        \includegraphics[width=0.5\textwidth]{Solution/AMC_10_10B_Fall_2021_11_Solution_Figure_1.png}
                    \end{figure}
                    

 This is the graph of the original Hexagon. After reflecting each minor arc over the sides of the hexagon it will look like this;

 
                    \begin{figure}[H]
                        \centering
                        \includegraphics[width=0.5\textwidth]{Solution/AMC_10_10B_Fall_2021_11_Solution_Figure_2.png}
                    \end{figure}
                    

 This bounded region is the same as the area of the hexagon minus the area of each of the reflect arcs. From the first diagram, the area of each arc is the area of the $60^{\circ}$ sector minus the area of the equilateral triangle, so each arc has an area of $\frac{\pi r^2}{6} - \frac{s^2\sqrt{3}}{4} \implies \frac{\pi}{6} - \frac{\sqrt{3}}{4}$.

 There are 6 total arcs, so the total area of the arcs is $6\cdot (\frac{\pi}{6} - \frac{\sqrt{3}}{4}) = \pi - \frac{3\sqrt{3}}{2}$.

 The area of the hexagon is $6\cdot \frac{\sqrt{3}}{4} = \frac{3\sqrt{3}}{2}$, so the area of the bounded region is:$\frac{3\sqrt{3}}{2} - (\pi - \frac{3\sqrt{3}}{2}) = 3\sqrt{3} - \pi = \boxed{B}$




\textbf{Solution 2}

Let the hexagon described be of area $H$ and let the circle's area be $C$. Let the area we want to aim for be $A$. Thus, we have that $C-H=H-A$, or $A=2H-C$. By some formulas, $C=\pi{r}^2=\pi$ and $H=6\cdot\frac12\cdot1\cdot(\frac12\cdot\sqrt3)=\frac{3\sqrt3}2$. Thus, $A=3\sqrt3-\pi$ or $\boxed{(\textbf{B})}$.

 



\textbf{Solution 3}

Denote by $O$ the center of this circle.Hence, the radius of this circle is 1.Denote this hexagon as $ABCDEF$.

 We have $\angle AOB = 60^\circ$.Hence, the area of the region formed between segment $AB$ and the minor arc formed by $A$ and $B$, denoted as $M$, is
 \begin{align*} M & = \pi 1^2 \frac{60}{360} - \frac{\sqrt{3}}{4} 1^2 \\ & = \frac{\pi}{6} - \frac{\sqrt{3}}{4} . \end{align*}


 Therefore, the area of the region that this problem asks us to compute is
 \begin{align*} \pi 1^2 - 12 M & = 3 \sqrt{3} - \pi . \end{align*}


 Therefore, the answer is $\boxed{\textbf{(B) }3 \sqrt{3} - \pi}$.

 



\textbf{Solution 4}

The area of the desired shape is equal to the $A($whole hexagon$)-A($Six arcs$).$ 

 Since the arcs are reflected upon the sides of the hexagon, we can see that $A($Six arcs$)=A($Whole circle$)-A($Whole hexagon$)$. 

 The area of the circle is $\pi\cdot{1^2}$, since the shape has side length $1$ and is inscribed within the circle (so its diameter is $2$).

 Combining these two, we see that 
 \[A(\text{Desired shape})=2A(\text{Hexagon})-A(\text{Circle}) \rightarrow\]
 \[A(\text{Desired shape})=2A(\text{Hexagon})-\pi.\]From here, three solutions can progress: 

 



\textbf{Solution 4.1 (General polygons)}

The area of a regular polygon is equal to $\frac{p\cdot{a}}{2}$, where $p$ is the perimiter of the polygon, and $a$ is the length of its apothem.

 The apothem is the distance from the center of the polygon to the midpoint of two adjacent vertices. If we were to create an equilateral triangle, whose base is at the side of the polygon, and its two sides meeting at the center, its height would be the apothem. 

 In this case, the side length of the hexagon is $1$. We can now split this equlilateral into two congruent right triangles, who both have the apothem as a side. Each triangle has side lengths of $\frac{1}{2}$,half of the base, $1$, the radius, and the apothem. By the pythagorean theorem, $1^2=(\frac{1}{2})^2+a^2$, for the apothem $a$. Solving yields $a=\frac{\sqrt{3}}{2}.$ 

 Since $p=6$, our formula becomes 
 \[A=6(\frac{\sqrt{3}}{2})\cdot{\frac{1}{2}}=\frac{3\sqrt{3}}{2}.\]

 Reall that the area of the figure is $2A(\text{Hexagon})-\pi.$ $2A=3\sqrt{3},$ so the sought area is $3\sqrt{3}-\pi$ , so our answer is $\boxed{\textbf{(B) }3 \sqrt{3} - \pi}$

 Alternatively, we could find the length of the apothem by the formula $\frac{s}{2\text{tan}(\frac{180}{n})},$ where $s$ is the side length and $n$ is the number of sides.

 



\textbf{Solution 4.2 (Quick)}

The area of a regular hexagon is given by the formula $\frac{3\sqrt{3}}{2}\cdot{s}$, for the side length $s$, which in this case, is $1$, so the area of this hexagon is $\frac{3\sqrt{3}}{2}.$

 We seek $2A(\text{Hexagon})-\pi,$ which is $\boxed{\textbf{(B) }3 \sqrt{3} - \pi}$




\textbf{Solution 5}

Use the diagrams provided in Solution $1$.

 The problem wants the area of the region bounded by the $6$ reflected arcs, a.k.a the area of the hexagon minus the area of the arcs drawn in the first figure. The area of those arcs can be calculated by taking the total area of the circle minus the total area of the hexagon.

 In the end, we are looking for the value of the (area of the hexagon)-((area of the circle)-(area of the hexagon)).

 The area of a hexagon with side length $1$ is $\frac{3\sqrt{3}}{2}$. The area of a circle with radius $1$ is $\pi$.

 $\frac{3\sqrt{3}}{2}-(\pi-\frac{3\sqrt{3}}{2})=3\sqrt{3}-\pi$.

 The answer is $\boxed{\textbf{(B) }3\sqrt{3}-\pi}$.

 


	\item (\textit{AMC 10 10B Fall 2021}) Which of the following conditions is sufficient to guarantee that integers $x$, $y$, and $z$ satisfy the equation

 \[x(x-y)+y(y-z)+z(z-x) = 1?\]
$\textbf{(A)} \: x>y$ and $y=z$
$\textbf{(B)} \: x=y-1$ and $y=z-1$
$\textbf{(C)} \: x=z+1$ and $y=x+1$
$\textbf{(D)} \: x=z$ and $y-1=x$
$\textbf{(E)} \: x+y+z=1$

 \textbf{Answer: }D



\textbf{Solution 1 (Completing the Square)}

It is obvious $x$, $y$, and $z$ are symmetrical. We are going to solve the problem by Completing the Square.

 $x ^ 2 + y ^ 2 + z ^ 2 - xy - yz - zx = 1$
$2x ^ 2 + 2y ^ 2 + 2z ^ 2 - 2xy - 2yz - 2zx = 2$
$(x-y)^2 + (y-z)^2 + (z-x)^2 = 2$

 Because $x, y, z$ are integers, $(x-y)^2$, $(y-z)^2$, and $(z-x)^2$ can only equal $0, 1, 1$. So one variable must equal another, and the third variable is $1$ different from those $2$ equal variables. So the answer is $\boxed{D}$.

 



\textbf{Solution 2}

Plugging in every choice, we see that choice $\textbf{(D)}$ works. 

 
We have $y=x+1, z=x$, so 
 \[x(x-y)+y(y-z)+z(z-x)=x(x-(x+1))+(x+1)((x+1)-x)+x(x-x)=x(-1)+(x+1)(1)=1.\]Our answer is $\textbf{(D)}$.

 



\textbf{Solution 3 (Bash)}

Just plug in all these options one by one, and one sees that all but $D$ fails to satisfy the equation.

 For $D$, substitute $z=x$ and $y=x+1$:

 $LHS=x(x-(x+1))+(x+1)(x+1-x)+x(x-x)=(-x)+(x+1)=1=RHS$

 Hence the answer is $\boxed{\textbf{(D)}}.$




\textbf{Solution 4 (Strategy)}

Looking at the answer choices and the question, the simplest ones to plug in would be equalities because it would make one term of the equation become zero. We see that answer choices A and D have the simplest equalities in them. However, A has an inequality too, so it would be simpler to plug in D which has another equality. We see that $x=z$ and $y-1=x$ means the equation becomes $x(x-(x+1)) + (x+1)(x+1 - x)  = 1 \implies -x + x + 1 = 1 \implies 1 =  1$, which is always true, so the answer is $\boxed{D}$




\textbf{Solution 5 (Logic and some Guesswork)}

We know that $x$, $y$, and $z$ are symmetric, which means that answer choices B and C are the same, so it can't be either of them. Answer choice E doesn't make sense because it is only 1 equation, which probably won't result in only 1 possible answer. This leaves us with choices A and D. Choice A isn't right because if $x$ is really big and $y=z=0$, then the equation becomes $x^2-x=1$ which doesn't work with a large integer $x$. Therefore, the only remaining answer choice is $\boxed{D}$



	\item (\textit{AMC 10 10B Fall 2021}) A square with side length $3$ is inscribed in an isosceles triangle with one side of the square along the base of the triangle. A square with side length $2$ has two vertices on the other square and the other two on sides of the triangle, as shown. What is the area of the triangle?


 
                    \begin{figure}[H]
                        \centering
                        \includegraphics[width=0.5\textwidth]{Problem/AMC_10_10B_Fall_2021_13_Problem_Figure_1.png}
                    \end{figure}
                    
$(\textbf{A})\: 19\frac14\qquad(\textbf{B}) \: 20\frac14\qquad(\textbf{C}) \: 21 \frac34\qquad(\textbf{D}) \: 22\frac12\qquad(\textbf{E}) \: 23\frac34$

 \textbf{Answer: }B



\textbf{Solution 1}

Let's split the triangle down the middle and label it:

 
                    \begin{figure}[H]
                        \centering
                        \includegraphics[width=0.5\textwidth]{Solution/AMC_10_10B_Fall_2021_13_Solution_Figure_1.png}
                    \end{figure}
                    

 We see that $\bigtriangleup ADG \sim \bigtriangleup BEG \sim \bigtriangleup CFG$ by AA similarity. $BE =  \frac{3}{2}$ because $AG$ cuts the side length of the square in half; similarly, $CF = 1$. Let $CG = h$: then by side ratios,

 
 \[\frac{h+2}{h} =  \frac{\frac{3}{2}}{1} \implies 2(h+2) = 3h \implies h = 4\].

 Now the height of the triangle is $AG = 4+2+3 =  9$. By side ratios, 
 \[\frac{9}{4} = \frac{AD}{1} \implies AD = \frac{9}{4}\]. 

 The area of the triangle is $AG\cdot AD = 9 \cdot \frac{9}{4}  = \frac{81}{4} = \boxed{\textbf{(B) }20 \frac{1}{4}}$




\textbf{Solution 2}

By similarity, the height is $3+\frac31\cdot2=9$ and the base is $\frac92\cdot1=4.5$. Thus the area is $\frac{9\cdot4.5}2=20.25=\boxed{\textbf{(B) }20 \frac{1}{4}}$.

 



\textbf{Solution 3 (With two different endings)}

This solution is based on this figure: \href{https://artofproblemsolving.com/wiki/index.php?title=File:2021\_AMC\_10B\_(Nov)\_Problem\_13,\_sol.png}{Image:2021\_AMC\_10B\_(Nov)\_Problem\_13,\_sol.png}

 Denote by $O$ the midpoint of $AB$.

 Because $FG = 3$, $JK = 2$, $FJ = KG$, we have $FJ = \frac{1}{2}$.

 We observe $\triangle ADF \sim \triangle FJH$.Hence, $\frac{AD}{FJ} = \frac{FD}{HJ}$.Hence, $AD = \frac{3}{4}$.By symmetry, $BE = AD = \frac{3}{4}$.

 Therefore, $AB = AD + DE + BE = \frac{9}{2}$.

 Because $O$ is the midpoint of $AB$, $AO = \frac{9}{4}$.

 We observe $\triangle AOC \sim \triangle ADF$.Hence, $\frac{OC}{DF} = \frac{AO}{AD}$.Hence, $OC = 9$.

 Therefore, ${\textrm Area} \ \triangle ABC = \frac{1}{2} AB \cdot OC = \frac{81}{4} = 20 \frac{1}{4}$.

 Therefore, the answer is $\boxed{\textbf{(B) }20 \frac{1}{4}}$.

 Alternatively, we can find the height in a slightly different way.

 Following from our finding that the base of the large triangle $AB = \frac{9}{2}$, we can label the length of the altitude of $\triangle{CHI}$ as $x$. Notice that $\triangle{CHI} \sim \triangle{CAB}$. Hence, $\frac{HI}{AB} = \frac{x}{CO}$. Substituting and simplifying, $\frac{HI}{AB} = \frac{x}{CO} \Rightarrow \frac{2}{\frac{9}{2}} = \frac{x}{x+5} \Rightarrow \frac{x}{x+5} = \frac{4}{9} \Rightarrow x = 4 \Rightarrow CO = 4 + 5 = 9$. Therefore, the area of the triangle is $\frac{\frac{9}{2} \cdot 9}{2} = \frac{81}{4} = \boxed{\textbf{(B) }20 \frac{1}{4}}$.

 



\textbf{Solution 4 (Coordinates)}

For convenience, we will use the image provided in the third solution.

 We can set $O$ as the origin.

 We know that $FG = 3$ and $JK = 2$. 

 We subtract $JK$ from $FG$ and divide by $2$ to get $KG = FJ = \frac{1}{2}$. 

 Since $HIKJ$ is a square, we know that $IK = 2$. 

 Using rise over run, we find that the slope of $CB$ is $\frac{-2}{0.5} = -4$.

 The coordinates of $I$ are $(1, 5)$. We plug this in to get the equation of the line that $CB$ runs along: 
 \[y = -4x + 9\] 

 We know that the $x-value$ of $C$ is $0$. Using this, we find that the $y-value$ is $9$. So the coordinates of $C$ are $(0, 9)$. 

 This gives us the height of $\triangle ACB$: $CO = 9$. 

 Now we need to find the coordinates of $B$.

 We know that the $y-value$ is $0$. Plugging this in, we find $0 = -4x +9$, or $\frac{9}{4} = x$.

 The coordinates of $B$ are $(\frac{9}{4}, 0)$.

 Since $\triangle ACB$ is symmetrical along $CO$, we can multiply $CO$ by $OB$ to get 
 \[9 \cdot \frac{9}{4} = \frac{81}{4}\]

 Simplifying, we get $\boxed{\textbf{(B) }20 \frac{1}{4}}$ for the area.

 



\textbf{Solution 5 (interesting)}

Call the right triangle to the right of the square with side length $2$ triangle $a_1$. Similarly, call the triangle to the right of the square with side length $3$ $a_2$. Label the ENTIRE triangle triangle ABC. We notice AA similarity between $a_1$ and $a_2$. Additionally, the base of $a_1$ must have a length of $0.5$, because the middle square takes up a length of $2$ and the other length of $1$ must be evenly split between the two congruent triangles next to the square, being all situated on a length of $3$. We then see that the ratio of $2:1$ must hold for the base of triangle ABC. $3:x = 2:1$, and we find that $x = 1.5$. The length of the base is then $1.5 + 3 = 9/2$.

 We can then find the height of the triangle by thinking about a geometrical sequence. It must be possible that you could continue inscribing squares in the triangles that are made, as every triangle created is similar to the one before it. Using this, we have the geometrical sequence of $3 + 2 + 4/3 +...$ (for each side length is the height that is being added) Using the geometrical sequence formula, we find $\frac{3}{1 - \frac{2}{3}}= \frac{3}{\frac{1}{3}} = 3\cdot3= 9$. The height is $9$ whilst the base length is $\frac{9}{2}$, and multiplying both values and dividing by $2$ yields $\frac{9(9)}{2(2)} = 81/4 = \boxed{\textbf{(B) }20 \frac{1}{4}}$



	\item (\textit{AMC 10 10B Fall 2021}) Una rolls $6$ standard $6$-sided dice simultaneously and calculates the product of the $6{ }$ numbers obtained. What is the probability that the product is divisible by $4?$
$\textbf{(A)}\: \frac34\qquad\textbf{(B)} \: \frac{57}{64}\qquad\textbf{(C)} \: \frac{59}{64}\qquad\textbf{(D)} \: \frac{187}{192}\qquad\textbf{(E)} \: \frac{63}{64}$

 \textbf{Answer: }C



\textbf{Solution}

We will use complementary counting to find the probability that the product is not divisible by $4$. Then, we can find the probability that we want by subtracting this from 1. We split this into $2$ cases. 

 Case 1: The product is not divisible by $2$.

 We need every number to be odd, and since the chance we roll an odd number is $\frac12,$ our probability is $\left(\frac12\right)^6=\frac1{64}.$

 Case 2: The product is divisible by $2$, but not by $4$.

 We need $5$ numbers to be odd, and one to be divisible by $2$, but not by $4$. There is a $\frac12$ chance that an odd number is rolled, a $\frac13$ chance that we roll a number satisfying the second condition (only $2$ and $6$ work), and $6$ ways to choose the order in which the even number appears.

 Our probability is $\left(\frac12\right)^5\left(\frac13\right)\cdot6=\frac1{16}.$

 Therefore, the probability the product is not divisible by $4$ is $\frac1{64}+\frac1{16}=\frac{5}{64}$.

 Our answer is $1-\frac{5}{64}=\boxed{\textbf{(C)}\ \frac{59}{64}}$.

 



\textbf{Solution 2(Pretty much same as above)}

Complementary counting implies that we can have either 0 factors of 2 or just 1 factor of 2.

 Case 1: 0 factors of 2

 The only odd numbers are 1, 3, 5 on a die so we have these 3 odd numbers with 6 positions so $3^6$ ways here.

 Case 2: 1 factor of 2

 The only even numbers that have only one factor of 2 on a die are 2, 6. On the 6 positions, there are 6 ways to place one of these even numbers and the rest of the 5 positions must be one of the three odd numbers. Since we could have either 2 or 6, we have

 $2 \cdot 6 \cdot 3^5 = 4 \cdot 3^6$. 

 Thus, we have $3^6 + 4 \cdot 3^6 = 5 \cdot 3^6$ cases that don't work.

 By complementary counting we get

 $1 - \frac{5 \cdot 3^6}{6^6} = 1 - \frac{5}{2^6} = \boxed{\frac{59}{64}}$



	\item (\textit{AMC 10 10B Fall 2021}) In square $ABCD$, points $P$ and $Q$ lie on $\overline{AD}$ and $\overline{AB}$, respectively. Segments $\overline{BP}$ and $\overline{CQ}$ intersect at right angles at $R$, with $BR=6$ and $PR=7$. What is the area of the square?


 
                    \begin{figure}[H]
                        \centering
                        \includegraphics[width=0.5\textwidth]{Problem/AMC_10_10B_Fall_2021_15_Problem_Figure_1.png}
                    \end{figure}
                    
$(\textbf{A})\: 85\qquad(\textbf{B}) \: 93\qquad(\textbf{C}) \: 100\qquad(\textbf{D}) \: 117\qquad(\textbf{E}) \: 125$

 \textbf{Answer: }D



\textbf{Solution 1}

Note that $\triangle APB \cong \triangle BQC$ by ASA. ($\angle PAB = \angle QBC = 90^\circ, AB=CB,$ and $\angle PBA = \angle QCB.$) Then, it follows that $\overline{PB} \cong \overline{QC}.$ Thus, $QC = PB = PR + RB = 7 + 6 = 13.$ Define $x$ to be the length of side $CR,$ then $RQ = 13-x.$ Because $\overline{BR}$ is the altitude of the triangle, we can use the property that $QR \cdot RC = BR^2.$ Substituting the given lengths, we have 
 \[(13-x) \cdot x = 36.\] Solving, gives $RQ = 4$ and $RC = 9.$ We eliminate the possibility of $x=4$ because $RC > QR.$ Thus, the side length of the square, by Pythagorean Theorem, is 
 \[\sqrt{9^2 +6^2} = \sqrt{81+36} = \sqrt{117}.\] Thus, the area of the square is $(\sqrt{117})^2 = 117,$ so the answer is $\boxed{\textbf{(D) }117}.$ 

 Please note that there is another way to prove that $CR = 4$ is impossible. If $CR = 4,$ then the side length would be $\sqrt{4^2 + 6^2} = \sqrt{52},$ and the area would be $52,$ but that isn't in the answer choices. Thus, $CR$ must be $9.$

 Extra Note: Another way to prove $4$ is impossible. The side length of the square, $S$, is equal to $\sqrt{4^2 + 6^2} = \sqrt{52}$. Because $x = 4$, $RQ = 9$. Because $QB = \sqrt{RB^2 + RQ^2} = \sqrt{6^2 + 9^2} = \sqrt{117}$ and $QB < S$ but $\sqrt{117} > \sqrt{52}$, we have proof by contradiction. And so $x = 9$.

 



\textbf{Solution 2 (Similarity, Pythagorean Theorem, and Systems of Equations)}

As above, note that $\bigtriangleup BPA \cong \bigtriangleup CQB$, which means that $QC =  13$. In addition, note that $BR$ is the altitude of a right triangle to its hypotenuse, so $\bigtriangleup BQR \sim \bigtriangleup CBR \sim \bigtriangleup CQB$. Let the side length of the square be $x$; using similarity side ratios of $\bigtriangleup BQR$ to $\bigtriangleup CQB$, we get 
 \[\frac{6}{x} = \frac{QB}{13} \implies QB \cdot x =  78\]Note that $QB^2 + x^2 = 13^2 = 169$ by the Pythagorean theorem, so we can use the expansion $(a+b)^2 = a^2+2ab+b^2$ to produce two equations and two variables;

 
 \[(QB + x)^2 = QB^2 + 2QB\cdot x + x^2 \implies (QB+x)^2 = 169 + 2 \cdot 78 \implies QB+x = \sqrt{13(13)+13(12)} = \sqrt{13 \cdot 25} = 5\sqrt{13}\]
 \[(QB-x)^2 = QB^2 - 2QB\cdot x + x^2 \implies (QB - x)^2 = 169 - 2\cdot 78 \implies \pm(QB-x) = \sqrt{13(13) - 13(12)}\]

 Since $QB-x$ is negative, it doesn't make sense in the context of this problem, so we go with 
 \[x-QB = \sqrt{13(13) - 13(12)} = \sqrt{13 \cdot 1} = \sqrt{13}\]

 We want $x^2$, so we want to find $x$. Adding the first equation to the second, we get 
 \[2x = 6\sqrt{13} \implies x = 3\sqrt{13}\]

 Then $x^2 = (3\sqrt{13}^2) = 9 \cdot 13 = 117 = \boxed{D}$

 -yingkai\_0\_ (Minor Edits)

 



\textbf{Solution 3}

We have that $\triangle CRB \sim \triangle BAP.$ Thus, $\frac{\overline{CB}}{\overline{CR}} = \frac{\overline{PB}}{\overline{AB}}$. Now, let the side length of the square be $s.$ Then, by the Pythagorean theorem, $CR = \sqrt{x^2-36}.$ Plugging all of this information in, we get 
 \[\frac{s}{\sqrt{s^2-36}} = \frac{13}{s}.\] Simplifying gives 
 \[s^2=13\sqrt{s^2-36},\] Squaring both sides gives 
 \[s^4 = 169s^2- 169\cdot 36 \implies s^4-169s^2 + 169\cdot 36 = 0.\] We now set $s^2=t,$ and get the equation $t^2-169t + 169\cdot 36 = 0.$ From here, notice we want to solve for $t$, as it is precisely $s^2,$ or the area of the square. So we use the \href{https://artofproblemsolving.com/wiki/index.php?title=Quadratic\_formula}{Quadratic formula}, and though it may seem bashy, we hope for a nice cancellation of terms. 
 \[t = \frac{169\pm\sqrt{169^2-4\cdot 36 \cdot 169}}{2}.\] It seems scary, but factoring $169$ from the square root gives us 
 \[t = \frac{169\pm \sqrt{169 \cdot (169-144)}}{2} = \frac{169 \pm \sqrt{169 \cdot 25}}{2} = \frac{169 \pm 13\cdot 5}{2} = \frac{169\pm 65}{2},\] giving us the solutions 
 \[t=52, 117.\] We instantly see that $t=52$ is way too small to be an area of this square ($52$ isn't even an answer choice, so you can skip this step if out of time) because then the side length would be $2\sqrt{13}$ and then, even the largest line you can draw inside the square (the diagonal)  is $2\sqrt{26},$ which is less than $13$ (line $PB$) And thus, $t$ must be $117$, and our answer is $\boxed{\textbf{(D)}}.$ $\blacksquare$




\textbf{Solution 4 (Point-line distance formula)}


                    \begin{figure}[H]
                        \centering
                        \includegraphics[width=0.5\textwidth]{Solution/AMC_10_10B_Fall_2021_15_Solution_Figure_1.png}
                    \end{figure}
                    Denote $a = RC$. Now tilt your head to the right and view $R, \overrightarrow{RB}$ and $\overrightarrow{RC}$ as the origin, $x$-axis and $y$-axis, respectively. In particular, we have points $B(6,0), C(0,a), P(-7,0)$. Note that side length of the square $ABCD$ is $BC = \sqrt{a^2 + 36}$. Also equation of line $BC$ is
 \[\underbrace{\frac{x}{6} + \frac{y}{a} = 1}_{\text{intercepts form}}     \quad \implies \quad     ax + 6y - 6a = 0.\]Because the distance from $P(-7,0)$ to line $\color[rgb]{0,0.4,0.65}BC: ax + 6y - 6a = 0$ is also the side length $\sqrt{a^2 + 36}$, we can apply the point-line distance formula to get
 \[\frac{|a\cdot(-7) + 6 \cdot 0 - 6a|}{\sqrt{a^2 + 36}} = {\sqrt{a^2 + 36}}\]which reduces to $|13a| = a^2 + 36$. Since $a$ is positive, the last equations factors as$a^2 - 13a + 36 = (a-4)(a-9) = 0$. Now judging from the figure, we learn that $a > RB = 6$. So $a = 9$.Therefore, the area of the square $ABCD$ is $BC^2 = RC^2 + RB^2 = a^2 + 6^2 = \boxed{\textbf{(D)}\ 117}$.

 



\textbf{Solution 5}

Denote $\angle PBA = \alpha$.Because $\angle QRB = \angle QBC = 90^\circ$, $\angle BCQ = \alpha$.

 Hence, $AB = BP \cos \angle PBA = 13 \cos \alpha$, $BC = \frac{BR}{\sin \angle BCQ} = \frac{6}{\sin \alpha}$.

 Because $ABCD$ is a square, $AB = BC$.Hence, $13 \cos \alpha = \frac{6}{\sin \alpha}$.

 Therefore,
 \begin{align*} \sin 2 \alpha & = 2 \sin \alpha \cos \alpha \\ & = \frac{12}{13} . \end{align*}


 Thus, $\cos 2 \alpha = \pm \frac{5}{13}$.

 $\textbf{Case 1}$: $\cos 2 \alpha = \frac{5}{13}$.

 Thus, $\cos \alpha = \sqrt{\frac{1 + \cos 2 \alpha}{2}} = \frac{3}{\sqrt{13}}$.

 Hence, $AB = 13 \cos \alpha = 3 \sqrt{13}$.

 Therefore, ${\textrm Area} \ ABCD = AB^2 = 117$.

 $\textbf{Case 2}$: $\cos 2 \alpha = - \frac{5}{13}$.

 Thus, $\cos \alpha = \sqrt{\frac{1 + \cos 2 \alpha}{2}} = \frac{2}{\sqrt{13}}$.

 Hence, $AB = 13 \cos \alpha = 2 \sqrt{13}$.

 However, we observe $BQ = \frac{BR}{\cos \alpha} = 3 \sqrt{13} > AB$.Therefore, in this case, point $Q$ is not on the segment $AB$.

 Therefore, this case is infeasible.

 Putting all cases together, the answer is $\boxed{\textbf{(D) }117}$.

 



\textbf{Solution 6 (Answer choices and areas)}

Note that if we connect points $P$ and $C$, we get a triangle with height $RC$ and length $13$. This triangle has an area of $\frac {1}{2}$ the square. We can now use answer choices to our advantage!

 Answer choice A: If $BC$ was $\sqrt {85}$, $RC$ would be $7$. The triangle would therefore have an area of $\frac {91}{2}$ which is not  half of the area of the square. Therefore, A is wrong.

 Answer choice B: If $BC$ was $\sqrt {93}$, $RC$ would be $\sqrt {57}$. This is obviously wrong.

 Answer choice C: If $BC$ was $10$, we would have that $RC$ is $8$. The area of the triangle would be $52$, which is not half the area of the square. Therefore, C is wrong.

 Answer choice D: If $BC$ was $\sqrt {117}$, that would mean that $RC$ is $9$. The area of the triangle would therefore be $\frac {117}{2}$ which IS half the area of the square. Therefore, our answer is $\boxed {\textbf{(D) 117}}$.

 



\textbf{Solution 7 (Power of a Point)}

Note that $PRQA$ is a cyclic quadrilateral (opposite angles add to $180^{\circ}$). Call the circumcircle of quadrilateral $PRQA$ $O$. Then the power of $B$ to $O$ is $6\cdot (6+7)=78$. Let $a$ be the length of $BQ$ and $s$ the side length of the square, then we have $a\cdot s = 78$, and we also have $a^2 + s^2=13^2$ (see note), solving the two equations will give us $s^2=117$.

 


 
            \begin{itemize}
        \item Note: This is from the fact that $\triangle QCB$ and $\triangle PBA$ are congruent. A short proof: We can prove that $\triangle PBA$ and $\triangle QBR$ are similar by ASA similarity. Likewise, we can prove that $\triangle QCB$ and $\triangle QBR$ are similar. Therefore, $\triangle QCB$ and $\triangle PBA$ are similar. Since they both have the side of a square as a hypotenuse, their hypotenuses are equal, so they are congruent. Therefore, $\overline{CQ}=\overline{PB}=13$.
            \end{itemize}
            



\textbf{Solution 8 (Pythagorean Theorem)}

Let's start by connecting the line $CP$ and calling the side length of the square $x$. 

 By the Pythagorean Theorem, we will have $CR^2+RB^2=CB^2$, so $CR^2+6^2=x^2$, so $CR=\sqrt{x^2-36}$. 

 Now, let's find $CP$ by using the Pythagorean theorem with the equation $CP^2=PR^2+CR^2$. Because $PR=7$ and $CR=\sqrt{x^2-36}$, we have $CP^2=7^2+(\sqrt{x^2-36})^2$, so $CP=\sqrt{x^2+13}$. Because $CD$ is a side of square $ABCD$, $CD=x$, and so now we can have $DP^2+CD^2=CP^2$, which will give us $DP=\sqrt{13}$. 

 Because $AB=x$, $AP=AD-DP=x-\sqrt{13}$, and $PB=6+7=13$, we now have $AB^2+AP^2=PB^2$ by the Pythagorean theorem, and the equation will be $x^2+(x-\sqrt{13})^2=169$. After solving this equation for $x$, we will get $x=-2\sqrt{13}$ or $3\sqrt{13}$, and because side lengths can't be negative, we will take the solution $x=3\sqrt{13}$, which means the area of the square will be $(3\sqrt{13})^2=117$. 

 Therefore, our answer is $\boxed{\textbf{(D) }117}.$



	\item (\textit{AMC 10 10B Fall 2021}) Five balls are arranged around a circle. Chris chooses two adjacent balls at random and interchanges them. Then Silva does the same, with her choice of adjacent balls to interchange being independent of Chris's. What is the expected number of balls that occupy their original positions after these two successive transpositions?


 $(\textbf{A})\: 1.6\qquad(\textbf{B}) \: 1.8\qquad(\textbf{C}) \: 2.0\qquad(\textbf{D}) \: 2.2\qquad(\textbf{E}) \: 2.4$

 \textbf{Answer: }D



\textbf{Solution 1}

After the first swap, we do casework on the next swap.

 
Case 1: Silva swaps the two balls that were just swapped

 There is only one way for Silva to do this, and it leaves 5 balls occupying their original position.

 
Case 2: Silva swaps one ball that has just been swapped with one that hasn't been swapped

 There are two ways for Silva to do this, and it leaves 2 balls occupying their original positions.

 
Case 3: Silva swaps two balls that have not been swapped

 There are two ways for Silva to do this, and it leaves 1 ball occupying their original position.

 
Our answer is the average of all 5 possible swaps, so we get 
 \[\frac{5+2\cdot2+2\cdot1}{5}=\frac{11}5=\boxed{(\textbf{D}) \: 2.2}.\]




\textbf{Solution 2 (Linearity of Expectation)}

The "expected value" in the question tips us off to this technique. Consider any ball. The probability it returns to the same position is the probability of being swapped twice plus the probability of never being swapped: 
 \[\frac{2}{5} \cdot \frac{1}{5} + \left(\frac{3}{5}\right)^2 = \frac{11}{25}.\] Multiply by 5 for 5 balls to get $\boxed{(\textbf{D}) \: 2.2}.$




\textbf{Solution 3 (Calculating-Mean)}

In probability and statistics, you learn that the mean and expected value are fundamentally equivalent values.

 The total number of ways to pick the balls is ${5 \choose 2}^2 = 100$.

 Now there are three cases either Chris and Silva choose the same two balls (C1), they choose two different balls (C2), or they choose one overlapping ball and three nonoverlapping (C3).

 C1: There are ${5 \choose 2} = 10$ possible placements of choosing $2$ balls. This gives a remaining $5$ unarranged balls. This gives a weighted mean pair of $10(5) = 50$.

 C2: There are $5$ ways of choosing $4$ balls adjacently, and due to Chris and Silva each choosing $2$ gives $10$ ways due to symmetry. This gives a weighted mean pair of $10(1) = 10$.

 C3: There are now $100 - 10 - 10 = 80$ ways for C3 to occur. This gives $2$ balls that remain in their original positions for a weighted mean pair of $80(2) = 160$.

 Calculating the mean gives $\frac{50 + 10 + 160}{100} = \boxed{(\textbf{D}) \: 2.2}$.

 


	\item (\textit{AMC 10 10B Fall 2021}) Distinct lines $\ell$ and $m$ lie in the $xy$-plane. They intersect at the origin. Point $P(-1, 4)$ is reflected about line $\ell$ to point $P'$, and then $P'$ is reflected about line $m$ to point $P''$. The equation of line $\ell$ is $5x - y = 0$, and the coordinates of $P''$ are $(4,1)$. What is the equation of line $m?$
$(\textbf{A})\: 5x+2y=0\qquad(\textbf{B}) \: 3x+2y=0\qquad(\textbf{C}) \: x-3y=0$
$(\textbf{D}) \: 2x-3y=0\qquad(\textbf{E}) \: 5x-3y=0$

 \textbf{Answer: }D



\textbf{Solution 1}

Denote $O$ as the origin.

 Even though the problem is phrased as a coordinate bash, that looks disgusting. Instead, let's try to phrase this problem in terms of Euclidean geometry, using the observation that $\angle POP'' = 90^{\circ}$, and that both $\ell$ and $m$ must pass through $O$ in order to preserve the distance from $P$ to the origin.
                    \begin{figure}[H]
                        \centering
                        \includegraphics[width=0.5\textwidth]{Solution/AMC_10_10B_Fall_2021_17_Solution_Figure_1.png}
                    \end{figure}
                    ($A$ and $B$ are just defined as points on lines $\ell$ and $m$.)Because of how reflections work, we have that $\angle AOP' = \angle POA$ and $\angle P'OB = \angle BOP''$; adding these two equations together and using angle addition, we have that $\angle AOB = \angle POA + \angle BOP''$. Since the sum of both sides combined must be $90^{\circ}$ by angle addition,
 \[\angle AOB = 45^{\circ}.\]This is helpful! We can now return to using coordinates, with this piece of information in mind:
                    \begin{figure}[H]
                        \centering
                        \includegraphics[width=0.5\textwidth]{Solution/AMC_10_10B_Fall_2021_17_Solution_Figure_2.png}
                    \end{figure}
                    The $45^{\circ}$ angle is a little bit unwieldy in the coordinate plane, so we should try to make a $45-45-90$ triangle. Let $A$ be a point on $\ell$; to make $A$ fit nicely in the diagram, let it be $(0.5,2.5)$. Now, let's draw a perpendicular to $\ell$ through point $A$, intersecting $m$ at point $B$. $OAB$ is a $45-45-90$ triangle, so $B$ is a $90$ degree counterclockwise rotation from $O$ about $A$. Therefore, the coordinates of $B$ are
 \[(0.5+2.5,2.5-0.5) = (3,2).\]So, $(3,2)$ is a point on line $m$, which we already know passes through the origin; therefore, $m$'s equation is $y=\frac{2x}{3} \implies \boxed{\textbf{(D) } 2x-3y = 0}.$

 (We never actually had to use the information of the exact coordinates of $P$; as long as $\angle POP'' = 90^{\circ}$, when we move $P$ around, this will not affect $m$'s equation.)

 \textbf{Supplement} In case you are confused about the coordinates of B, first transform O, A, and B such that A is the origin A'(0,0) in a new coordinate system. From there it is not too hard to see that O now has the coordinates O' (-0.5,-2.5). Thus point B, from a 90 deg CCW rotation around origin A, will have a coordinate of B' (2.5, -0.5). 

 Now in order to go from this new coord system A'(0,0) to its original point A(0.5,2.5), we have to +0.5,+2.5 to the x and y coordinates respectively. Doing this with B' to B, we have (2.5+0.5,-0.5+2.5)=B (3,2) as desired.

 



\textbf{Solution 2}

It is well known that the composition of 2 reflections , one after another, about two lines $l$ and $m$, respectively, that meet at an angle $\theta$ is a rotation by $2\theta$ around the intersection of $l$ and $m$. 

 Now, we note that $(4,1)$ is a 90 degree rotation clockwise of $(-1,4)$ about the origin, which is also where $l$ and $m$ intersect. So $m$ is a 45 degree rotation of $l$ about the origin clockwise. 

 To rotate $l$ 90 degrees clockwise, we build a square with adjacent vertices $(0,0)$ and $(1,5)$. The other two vertices are at $(5,-1)$ and $(6,4)$. The center of the square is at $(3,2)$, which is the midpoint of $(1,5)$ and $(5,-1)$. The line $m$ passes through the origin and the center of the square we built, namely at $(0,0)$ and $(3,2)$. Thus the line is $y = \frac{2}{3} x$. The answer is $\boxed{\textbf{(D) } 2x-3y = 0}$.

 



\textbf{Solution 3}

We know that the equation of line $\ell$ is $y = 5x$. This means that $P'$ is $(-1,4)$ reflected over the line $y = 5x$. This means that the line with $P$ and $P'$ is perpendicular to $\ell$, so it has slope $-\frac{1}{5}$. Then the equation of this perpendicular line is $y = -\frac{1}{5}x + c$, and plugging in $(-1,4)$ for $x$ and $y$ yields $c = \frac{19}{5}$. 

 The midpoint of $P'$ and $P$ lies at the intersection of $y = 5x$ and $y = -\frac{1}{5}x + \frac{19}{5}$. Solving, we get the x-value of the intersection is $\frac{19}{26}$ and the y-value is $\frac{95}{26}$. Let the x-value of $P'$ be $x'$ - then by the midpoint formula, $\frac{x' - 1}{2} = \frac{19}{26} \implies x' = \frac{32}{13}$. We can find the y-value of $P'$ the same way, so $P' = (\frac{32}{13},\frac{43}{13})$.

 Now we have to reflect $P'$ over $m$ to get to $(4,1)$. The midpoint of $P'$ and $P''$ will lie on $m$, and this midpoint is, by the midpoint formula, $(\frac{42}{13},\frac{28}{13})$. $y = mx$ must satisfy this point, so $m = \frac{\frac{28}{13}}{\frac{42}{13}} = \frac{28}{42} = \frac{2}{3}$.

 Now the equation of line $m$ is $y = \frac{2}{3}x \implies 2x-3y = 0 = \boxed{D}$




\textbf{Solution 4}

First, use Solution 1's method to get $\angle POP'' = 90^\circ$ and that the angle between lines $\ell$ and $m$ is $45^\circ$. From here, note that the slope of line $m$ is less than that of line $\ell$ as otherwise $P''$ wouldn't even be close to $(4, 1)$. Thus, line $m$ is a $45^\circ$ clockwise rotation of line $\ell$. Line $\ell$ makes an angle of $\tan^{-1}(5)$ with the positive x axis. Thus, line $m$ makes an angle of $\tan^{-1}(5) - 45^\circ$ with the positive x axis. Thus, the slope of line $m$ is
 \[\tan (\tan^{-1}(5) - 45^\circ) = \frac{5 - 1}{1 + 5\cdot 1} = \frac{2}{3},\]by the tangent addition formula. Since the slope of line $m$ is $\frac{2}{3}$, its equation is $y = \frac{2}{3}x \implies 2x - 3y = 0$, which is choice $\boxed{\textbf{D}}$.

 



\textbf{Solution 5(cheese)}

When we graph all the lines and points with a ruler, you can see that a slope of $\frac{5}{3}$ is too big while $\frac{1}{3}$ is too small. We also see that the slope cannot be negative, therefore the answer is $\boxed{D}.$~ agentdabber

 


	\item (\textit{AMC 10 10B Fall 2021}) Three identical square sheets of paper each with side length $6{ }$ are stacked on top of each other. The middle sheet is rotated clockwise $30^\circ$ about its center and the top sheet is rotated clockwise $60^\circ$ about its center, resulting in the $24$-sided polygon shown in the figure below. The area of this polygon can be expressed in the form $a-b\sqrt{c}$, where $a$, $b$, and $c$ are positive integers, and $c$ is not divisible by the square of any prime. What is $a+b+c?$

                    \begin{figure}[H]
                        \centering
                        \includegraphics[width=0.5\textwidth]{Problem/AMC_10_10B_Fall_2021_18_Problem_Figure_1.png}
                    \end{figure}
                    
$(\textbf{A})\: 75\qquad(\textbf{B}) \: 93\qquad(\textbf{C}) \: 96\qquad(\textbf{D}) \: 129\qquad(\textbf{E}) \: 147$

 \textbf{Answer: }E



\textbf{Solution 1 (Trigonometry)}


                    \begin{figure}[H]
                        \centering
                        \includegraphics[width=0.5\textwidth]{Solution/AMC_10_10B_Fall_2021_18_Solution_Figure_1.png}
                    \end{figure}
                    The $24$-sided polygon is made out of $24$ shapes like $\triangle ABC$. Then $\angle BAC=360^\circ/24=15^\circ$, and $\angle EAC = 45^\circ$, so  $\angle{EAB} = 30^{\circ}$. Then $EB=AE\tan 30^\circ = \sqrt{3}$; therefore $BC=EC-EB=3-\sqrt{3}$. Thus
 \[[ABC] = \frac{BC}{EC}\cdot [ACE] = \frac{3-\sqrt{3}}{3}\cdot \frac 92\]and the required area is $24\cdot[ABC] =108-36\sqrt{3}$. Finally $108+36+3=\boxed{(\textbf{E})\ 147}$. ~lopkiloinm

 



\textbf{Solution 2 (15-75-90 triangles)}


                    \begin{figure}[H]
                        \centering
                        \includegraphics[width=0.5\textwidth]{Solution/AMC_10_10B_Fall_2021_18_Solution_Figure_2.png}
                    \end{figure}
                    

 


 
Break the $24$-gon as shown so that there are $12$ copies of quadrilateral $OBAC$. We can find the area of this quadrilateral by finding the area of triangle $OBC$ and subtracting the area of triangle $ABC$.

 The angle from $O$ to one of the vertices of the original square is $45^\circ{}$, and this point rotates $30^\circ{}$, so the angle that $OC$ makes with the horizontal is $75^\circ{}$.

 Let $D$ be the intersection of $OA$ with $BC$. Since $ABC$ and $OBC$ are both isosceles, $OD$ forms a right angle with $BC$. Furthermore, $\angle DOC = 90^\circ{}-75^\circ{} = 15^\circ{}$.

 
                    \begin{figure}[H]
                        \centering
                        \includegraphics[width=0.5\textwidth]{Solution/AMC_10_10B_Fall_2021_18_Solution_Figure_3.png}
                    \end{figure}
                    

 Note that $OC$ is equal to $3\sqrt{2}$ because it is half the diagonal of the square with side length $6$. We can now split up the $15-75-90$ into $30-60-90$ and $15-15-150$ triangles to find the height and length of the triangle: 

 
                    \begin{figure}[H]
                        \centering
                        \includegraphics[width=0.5\textwidth]{Solution/AMC_10_10B_Fall_2021_18_Solution_Figure_4.png}
                    \end{figure}
                    

 Now we have use Pythagorean: $x^2+(x \cdot (2+\sqrt{3}))^2 = 3\sqrt{2}^2 \implies x^2+x^2 \cdot (7+4\sqrt{3}) = 18$$\implies x^2 = \frac{18}{8+4\sqrt{3}} = \frac{9}{4+2\sqrt{3}} \implies x^2 = \frac{36-18\sqrt{3}}{4}$

 Now note that $x$ must be in the form of $\frac{a\sqrt{3}-b}{2}$. Therefore, $-2ab = -18 \implies ab=9$. We can guess that $a=b=3$. Sure enough, $3\sqrt{3}^2 +3^2= 36$, so $x = \frac{3\sqrt{3}-3}{2}$. Then the height of the triangle is $x \cdot (2+\sqrt{3}) = \frac{3\sqrt{3}+3}{2}$.

 The area of triangle $OBC$ is equal to the length times the height of triangle $ODC$ since this triangle has half the area of the full triangle. Therefore the area of triangle $OBC = \frac{3\sqrt{3}-3}{2} \cdot \frac{3\sqrt{3}+3}{2} = \frac{18}{4}$.

 Now we look at triangle $ABC$. $\angle OCA = 45^\circ{}$ because $OC$ is the diagonal of a square, and $\angle COA$ is $15^\circ{}$ as we saw earlier. Therefore, $\angle OAC = \angle OAB = 120^\circ{}$. Because both these angles are $120$, $\angle BAC = 360-240 = 120$.

 We can now split triangle $ABC$ into two congruent $30-60-90$ triangles, and we know the base of each since $BC = 3\sqrt{3}-3$.

 The height of these triangles is $\frac{\frac{3\sqrt{3}-3}{2}}{\sqrt{3}} = \frac{3-\sqrt{3}}{2}$. Therefore the area of this triangle is $\frac{3\sqrt{3}-3}{2} \cdot \frac{3-\sqrt{3}}{2} = \frac{12\sqrt{3}-18}{4}$

 Therefore the total area of this segment is $\frac{18}{4}-(\frac{12\sqrt{3}-18}{4}) = \frac{36-12\sqrt{3}}{4} = 9-3\sqrt{3}$. Multiplying by $12$ to find the entire area, the area of the figure is $108-36\sqrt{3}$ and the answer is $108+36+3 = 147 = \boxed{E}$




\textbf{Solution 3}

First note the useful fact that if $R$ is the circumradius of a dodecagon ($12$-gon) the area of the figure is $3R^2.$ If we connect the vertices of the $3$ squares we get a dodecagon. The radius of circumcircle of the dodecagon is simply half the diagonal of the square, which is $3\sqrt{2}.$ Thus the area of the dodecagon is $3 \cdot (3\sqrt{2})^2 = 3 \cdot 18 = 54.$ But, the problem asked for the area of the combined figure which was made from the rotated squares. This area is the area of the dodecagon, which was found, subtracting the $12$ isosceles triangles, which are formed when connecting the vertices of the squares to created the dodecagon. To find this area, we need to know the base of the isosceles triangle, call this $x.$ Then, we can use the Law of Cosines on the triangle that is formed from the two vertices of the square and the center of the square. After computing, we get that $x = 3\sqrt{3} -3.$ Realize that the $12$ isosceles are congruent with an angle measure of $120^{\circ},$ this means that we can create $4$ congruent equilateral triangles with side length of $3\sqrt3 - 3.$ The area of the equilateral triangle is $\frac{\sqrt{3}}{4} \cdot (3\sqrt{3} -3)^2 =  \frac{\sqrt{3}}{4} \cdot (36 - 18\sqrt{3}) = \frac{36\sqrt{3} - 54}{4}.$ Thus, the area of all the twelve small equilateral traingles are $4 \cdot \frac{36\sqrt{3} - 54}{4} = 36\sqrt{3} - 54$. Thus, the requested area is $54 - (36\sqrt{3} - 54) = 108 - 36\sqrt{3}.$ Thus, $a+b+c = 108 + 36 + 3 = 147,$ so the answer is $\boxed{(\textbf{E})\textbf{147}}.$




\textbf{Solution 4 (30-60-90 Triangles)}

To make things simpler, let's take only the original sheet and the 30 degree rotated sheet. Then the diagram is this;

 
                    \begin{figure}[H]
                        \centering
                        \includegraphics[width=0.5\textwidth]{Solution/AMC_10_10B_Fall_2021_18_Solution_Figure_5.png}
                    \end{figure}
                    

 The area of this diagram is the original square plus the area of the four triangles that 'jut' out of the square. Because the square is rotated $30^{\circ}$, each triangle is a 30-60-90 triangle. Similarly, the triangles that are bounded on the inside of the original square outside of the rotated square are also congruent 30-60-90 triangles. Noting this, we can do some labelling:

 
                    \begin{figure}[H]
                        \centering
                        \includegraphics[width=0.5\textwidth]{Solution/AMC_10_10B_Fall_2021_18_Solution_Figure_6.png}
                    \end{figure}
                    

 Since the side lengths of the squares must be the same, and they are both 6, we have a system of equations;
 \[y+x+y\sqrt{3} = 6\]
 \[\frac{x\sqrt{3}}{2} + 2y + \frac{x}{2} = 6\]

 We solve this to get $x = 6-2\sqrt{3}$ and $y = 3-\sqrt{3}$.

 The area of each triangle is $\frac{x}{2} \cdot \frac{x\sqrt{3}}{2} \cdot \frac{1}{2}  = 6\sqrt{3} - 9$ by plugging in $x$.

 The rotated 60 degree square is the same thing as rotating it 30 degrees counterclockwise, so it's triangles that jut out of the square will be congruent to the triangles we have found, and therefore they will have the same area.

 Unfortunately, when drawing all three squares, we see the two triangles overlap; take the very top for example.

 
                    \begin{figure}[H]
                        \centering
                        \includegraphics[width=0.5\textwidth]{Solution/AMC_10_10B_Fall_2021_18_Solution_Figure_7.png}
                    \end{figure}
                    

 The area of this shape is twice the area of each of the triangles that we have already found minus the area of the small triangle that is overlapped by the two by PIE. Now we only need to find the area of $\bigtriangleup BCF$.

 $\angle GBD \cong \angle ECA \cong 30^{\circ}$ and by symmetry $\bigtriangleup BCF$ is isosceles, so it is a 30-30-120 triangle. If we draw a perpendicular, we split it into two 30-60-90 triangles;

 
                    \begin{figure}[H]
                        \centering
                        \includegraphics[width=0.5\textwidth]{Solution/AMC_10_10B_Fall_2021_18_Solution_Figure_8.png}
                    \end{figure}
                    

 By symmetry, the distance from A to the edge of the square is equal to the distance from D to the edge of the square is equal to $y$. AC = BD = $x$, and the side length of the square is 6, so we use PIE to obtain 
 \[x+x-BC = 6-y-y \implies BC = 12 - 6\sqrt{3}\]To find the height of $\bigtriangleup BFC$, we see that $HC = \frac{BC}{2} = 6-3\sqrt{3}$. Then by 30-60-90 triangles, $HF = \frac{HC}{\sqrt{3}} = 2\sqrt{3} - 3$. Finally, the area of $\bigtriangleup BFC = \frac{BC \cdot HF}{2} = 21\sqrt{3}-36$.

 Putting it all together, the area of the entire diagram is the area of the square plus four of these triangle-triangle intersections. The area of these intersections by PIE is $2 \cdot [ACE] - [BFC] = 12\sqrt{3}-18-(21\sqrt{3}-36) = 18-9\sqrt{3}$. Therefore the total area is $36 + 4 \cdot(18-9\sqrt{3}) = 36 + 72 - 36 \sqrt{3} = 108 - 36\sqrt{3}$.

 Thus $a + b + c = 108+36+3 = 147 = \boxed{\textbf{(E) }147}$.

 



\textbf{Solution 5}

As shown in \href{https://artofproblemsolving.com/wiki/index.php?title=File:2021\_AMC\_12B\_(Nov)\_Problem\_15,\_sol.png}{Image:2021\_AMC\_12B\_(Nov)\_Problem\_15,\_sol.png}, all 12 vertices of three squares form a regular dodecagon (12-gon).Denote by $O$ the center of this dodecagon.

 Hence, $\angle AOB = \frac{360^\circ}{12} = 30^\circ$.

 Because the length of a side of a square is 6, $AO = 3 \sqrt{2}$.

 Hence, $AB = 2 AO \sin \frac{\angle AOB}{2} = 3 \left( \sqrt{3} - 1 \right)$.

 We notice that $\angle MAB = \angle MBA = 30^\circ$.Hence, $AM = \frac{AB}{2\cos \angle MAB} = 3 - \sqrt{3}$.

 Therefore, the area of the region that three squares cover is
 \begin{align*} &  \ [ABCDEFGHIJKL] - 12[MAB] \\ & = 12 [OAB] - 12 [MAB] \\ & = 12 \cdot \frac{1}{2} OA \cdot OB \sin \angle AOB - 12 \cdot \frac{1}{2} MA \cdot MB \sin \angle AMB \\ & = 6 OA^2 \sin \angle AOB - 6 MA^2 \sin \angle AMB \\ & = 108 - 36\sqrt{3} \end{align*}


 Therefore, the answer is $\boxed{\textbf{(E) }147}$.

 Note: If you have square brackets around a specified polygon (i.e. [ABCD]), it means the writer is referring to the area of that polygon.

 



\textbf{Solution 6}


                    \begin{figure}[H]
                        \centering
                        \includegraphics[width=0.5\textwidth]{Solution/AMC_10_10B_Fall_2021_18_Solution_Figure_9.png}
                    \end{figure}
                    First, we can separate the shape into 12 congruent kites.  The area of the figure can be determined by finding the area of one kite and multiplying it by 12.  In order to to get the area of one kite, we need to find its diagonals, shown in blue.  

 
                    \begin{figure}[H]
                        \centering
                        \includegraphics[width=0.5\textwidth]{Solution/AMC_10_10B_Fall_2021_18_Solution_Figure_10.png}
                    \end{figure}
                    We notice that angle FCE is $\frac{360^\circ}{12} = 30^\circ$. Also, we know that CD is half of AB, so it has a length of 3.  Now, we can find the lengths of FC and FD using the 30-60-90 triangle.  We find that FC is $2\sqrt{3}$ and FD is $\sqrt{3}$.  Since FC is congruent to CE, CE is also $2\sqrt{3}$.  Using this information, we can conclude that ED is $2\sqrt{3}-3$.

 
                    \begin{figure}[H]
                        \centering
                        \includegraphics[width=0.5\textwidth]{Solution/AMC_10_10B_Fall_2021_18_Solution_Figure_11.png}
                    \end{figure}
                    Now, we can find the shorter diagonal by using the Pythagorean theorem:
 \begin{align*} & {\textrm FC}^2 =  \sqrt{3}^2 + (2\sqrt{3}-3)^2 \\ & {\textrm FC} = \sqrt{24-  12\sqrt{3}} \end{align*}


                    \begin{figure}[H]
                        \centering
                        \includegraphics[width=0.5\textwidth]{Solution/AMC_10_10B_Fall_2021_18_Solution_Figure_12.png}
                    \end{figure}
                    We can find the longer diagonal of the kite by looking at one of the square sheets of paper.  We know that the side of the square has a length of 6, so the diagonal of the square must be $6\sqrt{2}$.  The longer diagonal of the kite is half of this length, so it has a length of $3\sqrt{2}$.

 The area of the entire figure is 
 \begin{align*} & = 12 \cdot\frac{{\textrm d_1*d_2}}{2} \\ & = 12 \cdot\frac{3\sqrt{2}\cdot\sqrt{24-  12\sqrt{3}}}{2} \\ & = 12 \cdot\frac{6\sqrt{12-  6\sqrt{3}}}{2} \\ & = 36 \sqrt{12-  6\sqrt{3}} \end{align*}


 Now we can use algebra to make our answer look a little nicer.
 \begin{align*} & a-\sqrt{b} = \sqrt{12-  6\sqrt{3}} \\ & (a-\sqrt{b})^2 = (\sqrt{12-  6\sqrt{3}})^2 \\ & a^2-2a\sqrt{b}+b = 12 - 6\sqrt{3}\\ &\\ &a^2+b = 12 \\ &2a\sqrt{b} = 6\sqrt{3} \\ & a = 3, b = 3 \\ &\\ &  a-\sqrt{b} = \sqrt{12-  6\sqrt{3}} = 3 - \sqrt{3} \end{align*}


 The area of the entire region is $36(3 - \sqrt{3})$, or $108 - 36\sqrt{3}$.  

 Therefore, $a + b + c = 108 + 36+ 3 = 147 = \boxed{E}$.

 



\textbf{Solution 7}

Let $O$ be the center of the polygon, $A$ be the bottom right corner of the first square, $C$ be the next vertex to the left of $A$, and $M$ be the midpoint between $A$ and $B$, where $B$ is the bottom left corner of the first square. Note that because there are three $90^{\circ}$ squares separated by $\frac{90^{\circ}}{3} = 30^{\circ}$, each side of the 24-sided polygon is equal in length, meaning to calculate the area of the whole polygon, we find the area of $\bigtriangleup OAC$ and multiply by 24.

 
To find $[\bigtriangleup OAC]$, we already know the height $\overline{OM}$ is the sidelength of the square over $2$, or $\frac{6}{2}=3$, so we just need the length of the base $\overline{AC}$. Notice that $\bigtriangleup OCM$ is a $30-60-90$ triangle since $\angle COM = \frac{360^{\circ}}{12} = 30^{\circ}$, so $\overline{CM} = \frac{\overline{OM}}{\sqrt{3}} = \frac{3}{\sqrt{3}} = \sqrt{3}$. Then $\overline{AC} = \overline{AM} - \overline{CM} = \frac{6}{2} - \sqrt{3} = 3 - \sqrt{3}$, so 
 \begin{align*} & [\bigtriangleup AOC] = \frac{1}{2} \cdot \overline{OM} \cdot \overline{AC} \\ & = \frac{1}{2} (3)(3 - \sqrt{3}) \\ & = \frac{9 - 3\sqrt{3}}{2} \end{align*}


 Then the whole area of the polygon is $\frac{9 - 3\sqrt{3}}{2} \cdot 24 = 108 - 36\sqrt{3}$. The desired solution is then $108 + 36 + 3 = 147$, so the answer is $\boxed{\textbf{(E) 147}}$.

 



\textbf{Solution 8}

Using the diagram from Solution 2, we observe two isosceles triangles. The larger triangle has angles of 75°, 75°, and 30°.

 To find the angle measures of the smaller triangle \(BCO\), note that the congruent sides of \(\triangle BCO\) are each half of the diagonals of two squares. Since a diagonal of a square bisects the right angle (45°), we have \(75 - 45 = 30^\circ\) for each base angle, making the remaining vertex angle \(120^\circ\).

 Because both triangles are isosceles, we have \(BO = CO\) and \(AB = AC\).

 Using the formula for the area of a square, $A = \frac{d^2}{2}$, we find

 $72 = \frac{d^2}{2} \Rightarrow d = 6\sqrt{2},$

 so $BO = 3\sqrt{2}.$

 Let $[BCO]$ denote the area of \(\triangle BCO\). Using the area formula $\frac{1}{2}ab\sin C$, we find

 $[BCO] = \frac{1}{2}(3\sqrt{2})(3\sqrt{2})\sin(30^\circ) = \frac{1}{2}(18)\left(\frac{1}{2}\right) = \frac{9}{2}.$

 Next, to find the side lengths of \(\triangle ABC\), we apply the Law of Cosines to \(\triangle BCO\):

 $BC^2 = (3\sqrt{2})^2 + (3\sqrt{2})^2 - 2(3\sqrt{2})(3\sqrt{2})\cos(120^\circ) = 36(1 - \frac{\sqrt{3}}{2}) = 18(2 - \sqrt{3}).$

 Now, let $AB = x.$ Applying the Law of Cosines to \(\triangle ABC\):

 $2x^2 - 2x^2\cos(120^\circ) = 18(2 - \sqrt{3}).$

 Simplifying,

 $3x^2 = 18(2 - \sqrt{3}) \Rightarrow x^2 = 6(2 - \sqrt{3}).$

 We can now find the area of \(\triangle ABC\) using the same area formula:

 $[ABC] = \frac{1}{2}x^2\sin(120^\circ) = \frac{1}{2}\bigl(6(2 - \sqrt{3})\bigr)\left(\frac{\sqrt{3}}{2}\right) = 3\sqrt{3} - \frac{9}{2}.$

 Subtracting the two areas gives the area of the quadrilateral \(OBAC\):

 $[OBAC] = 9 - 3\sqrt{3}.$

 Since there are 12 such regions in the figure, the total area is

 $12(9 - 3\sqrt{3}) = 108 - 36\sqrt{3}.$

 The problem asks for \(a + b + c,\) so

 $108 + 36 + 3 = \boxed{\textbf{(E) 147}}$



	\item (\textit{AMC 10 10B Fall 2021}) Let $N$ be the positive integer $7777\ldots777$, a $313$-digit number where each digit is a $7$. Let $f(r)$ be the leading digit of the $r{ }$th root of $N$. What is
 \[f(2) + f(3) + f(4) + f(5)+ f(6)?\]$(\textbf{A})\: 8\qquad(\textbf{B}) \: 9\qquad(\textbf{C}) \: 11\qquad(\textbf{D}) \: 22\qquad(\textbf{E}) \: 29$

 \textbf{Answer: }A



\textbf{Solution 1}

We can rewrite $N$ as $\frac{7}{9}\cdot 9999\ldots999 = \frac{7}{9}\cdot(10^{313}-1)$. When approximating values, as we will shortly do, the minus one will become negligible so we can ignore it.When we take the power of ten out of the square root, we’ll be multiplying by another power of ten, so the leading digit will not change. Thus the leading digit of $f(r)$ will be equal to the leading digit of $\sqrt[r]{\frac{7}{9} \cdot 10^{313(\mod r)}}$. 

 Then $f(2)$ is the first digit of $\sqrt{\frac{7}{9}\cdot(10)} = \sqrt{\frac{70}{9}} = \sqrt{7. \ldots} \approx 2$
$f(3) = \sqrt[3]{\frac{7}{9} \cdot 10} = \sqrt[3]{\frac{70}{9}} = \sqrt[3]{7.\ldots} \approx 1$.

 $f(4) = \sqrt[4]{\frac{7}{9} \cdot 10} = \sqrt[4]{\frac{70}{9}} = \sqrt[4]{7.\ldots} \approx 1$.

 $f(5) = \sqrt[5]{\frac{7}{9} \cdot 1000} = \sqrt[5]{\frac{7000}{9}} = \sqrt[5]{777.\ldots} \approx 3$.

 $f(6) = \sqrt[6]{\frac{7}{9} \cdot 10} = \sqrt[6]{\frac{70}{9}} = \sqrt[6]{7.\ldots} \approx 1$.

 The final answer is therefore $2+1+1+3+1 = \boxed{\textbf{(A) }8}$.

 Note: in all of the divisions, we omitted the decimal places, because they are irrelevant to finding the leading digit.

 



\textbf{Solution 2}

For notation purposes, let $x$ be the number $777 \ldots 777$ with $313$ digits, and let $B(n)$ be the leading digit of $n$. As an example, $B(x) = 7$, because $x = 777 \ldots 777$, and the first digit of that is $7$.

 Notice that 
 \[B(\sqrt{\frac{n}{100}}) = B(\sqrt{n})\] for all numbers $n \geq 100$; this is because $\sqrt{\frac{n}{100}} = \frac{\sqrt{n}}{10}$, and dividing by $10$ does not affect the leading digit of a number. Similarly, 
 \[B(\sqrt[3]{\frac{n}{1000}}) = B(\sqrt[3]{n}).\]In general, for positive integers $k$ and real numbers $n > 10^{k}$, it is true that
 \[B(\sqrt[k]{\frac{n}{10^{k}}}) = B(\sqrt[k]{n}).\]Behind all this complex notation, all that we're really saying is that the first digit of something like $\sqrt[3]{123456789}$ has the same first digit as $\sqrt[3]{123456.789}$ and $\sqrt[3]{123.456789}$.

 The problem asks for
 \[B(\sqrt[2]{x}) + B(\sqrt[3]{x}) + B(\sqrt[4]{x}) + B(\sqrt[5]{x}) + B(\sqrt[6]{x}).\]

 From our previous observation, we know that
 \[B(\sqrt[2]{x}) = B(\sqrt[2]{\frac{x}{100}} = B(\sqrt[2]{\frac{x}{10,000}} = B(\sqrt[2]{\frac{x}{1,000,000}} = \ldots .\] Therefore, $B(\sqrt[2]{x}) = B(\sqrt[2]{7.777 \dots})$. We can evaluate $B(\sqrt[2]{7.777 \dots})$, the leading digit of $\sqrt[2]{7.777 \dots}$, to be $2$. Therefore, $f(2) = 2$.

 Similarly, we have  
 \[B(\sqrt[3]{x}) = B(\sqrt[3]{\frac{x}{1,000}} = B(\sqrt[3]{\frac{x}{1,000,000}} = B(\sqrt[3]{\frac{x}{1,000,000,000}} = \ldots .\] Therefore, $B(\sqrt[3]{x}) = B(\sqrt[3]{7.777 \ldots})$. We know $B(\sqrt[3]{7.777 \ldots}) = 1$, so $f(3) = 1$.

 Next,
 \[B(\sqrt[4]{x}) = B(\sqrt[4]{7.777 \ldots})\]and $B(\sqrt[4]{7.777 \ldots}) = 1$, so $f(4) = 1$.

 We also have
 \[B(\sqrt[5]{x}) = B(\sqrt[5]{777.777 \ldots})\]and $B(\sqrt[5]{777.777 \ldots}) = 3$, so $f(5) = 3$.

 Finally,
 \[B(\sqrt[6]{x}) = B(\sqrt[6]{7.777 \ldots})\]and $B(\sqrt[4]{7.777 \ldots}) = 1$, so $f(6) = 1$.

 We have that $f(2)+f(3)+f(4)+f(5)+f(6) = 2+1+1+3+1 = \boxed{\textbf{(A) } 8}$.

 



\textbf{Solution 3 (Condensed Solution 1)}

Since $7777..7$ is a $313$ digit number and $\sqrt {7}$ is around $2.5$, we have $f(2)$ is $2$. $f(3)$ is the same story, so $f(3)$ is $1$. It is the same as $f(4)$ as well, so $f(4)$ is also $1$. However, $313$ is $3$ mod $5$, so we need to take the 5th root of $777$, which is between $3$ and $4$, and therefore, $f(5)$ is $3$. $f(6)$ is the same as $f(4)$, since it is $1$ more than a multiple of $6$. Therefore, we have $2+1+1+3+1$ which is $\boxed {\textbf{(A)8}}$.

 



\textbf{Solution 4}

First, we compute $f \left( 2 \right)$.

 Because $N > 4 \cdot 10^{312}$, $\sqrt{N} > 2 \cdot 10^{156}$.Because $N < 9 \cdot 10^{312}$, $\sqrt{N} < 3 \cdot 10^{156}$.

 Therefore, $f \left( 2 \right) = 2$.

 Second, we compute $f \left( 3 \right)$.

 Because $N > 1 \cdot 10^{312}$, $\sqrt[3]{N} > 1 \cdot 10^{104}$.Because $N < 8 \cdot 10^{312}$, $\sqrt[3]{N} < 2 \cdot 10^{104}$.

 Therefore, $f \left( 3 \right) = 1$.

 Third, we compute $f \left( 4 \right)$.

 Because $N > 1 \cdot 10^{312}$, $\sqrt[4]{N} > 1 \cdot 10^{78}$.Because $N < 16 \cdot 10^{312}$, $\sqrt[4]{N} < 2 \cdot 10^{78}$.

 Therefore, $f \left( 4 \right) = 1$.

 
Fourth, we compute $f \left( 5 \right)$.

 Because $N > 3^5 \cdot 10^{310}$, $\sqrt[5]{N} > 3 \cdot 10^{62}$.Because $N < 4^5 \cdot 10^{310}$, $\sqrt[5]{N} < 4 \cdot 10^{62}$.

 Therefore, $f \left( 5 \right) = 3$.

 Fifth, we compute $f \left( 6 \right)$.

 Because $N > 1 \cdot 10^{312}$, $\sqrt[6]{N} > 1 \cdot 10^{52}$.Because $N < 2^6 \cdot 10^{312}$, $\sqrt[6]{N} < 2 \cdot 10^{52}$.

 Therefore, $f \left( 6 \right) = 1$.

 Therefore,
 \begin{align*} f \left( 2 \right) + f \left( 3 \right) + f \left( 4 \right) + f \left( 5 \right) + f \left( 6 \right) & = 2 + 1 + 1 + 3 + 1 \\ & = 8 . \end{align*}


 Therefore, the answer is $\boxed{\textbf{(A) }8}$.

 



\textbf{Solution 5 (Guessing)}

Benford's Law states that in random numbers, the leading digit is more likely to be $1,2,$ or $3$ rather than $8$ or $9$. From here, we can eliminate C, D, E. It is better to guess between A and B than not guess at all since your expected score from doing this is $3$ points.

 



\textbf{Solution 6 (Intuition)}

To begin with, you can test some numbers and find a pattern. For $f(2)$, we get the leading digit of $\sqrt{7}$ is 2, and the leading digit of $\sqrt{77}$ is 8. With each $7$ afterward, it repeats in a loop: $2$ for $7$, $8$ for $77$, $2$ for $777$, $8$ for $777$, and so on. Since the problem states that $N$ is a $313$-digit number, we get the leading digit as $2$ for $f(2)$. For $f(3)$, it repeats in a loop of $3$ ($\text{3rd}$ root means it repeats every $3$ times). In this case, you would get the leading digit for the $\text{3rd}$ root of $7$ to be 1, $77$ it will be 4, and for $777$ it will be 9. Since $N$ is $1$ more than a multiple of 3, the leading digit for $f(3)$ is $1$. Likewise, we can find the leading digits for $f(4)$ to be $1$, $f(5)$ to be $3$, and $f(6)$ to be $1$:

 For $f(2)$: $7: \boxed{2},\hspace{0.085in} 77: 8 ...$

 For $f(3)$: $7: \boxed{1},\hspace{0.085in} 77: 4,\hspace{0.085in} 777: 9 ...$

 


 For $f(4)$: $7: \boxed{1},\hspace{0.085in} 77: 2,\hspace{0.085in} 777: 5,\hspace{0.085in} 7777: 9 ...$

 For $f(5)$: $7: 1,\hspace{0.085in} 77: 2,\hspace{0.085in} 777: \boxed{3},\hspace{0.085in} 7777: 6,\hspace{0.085in} 77777: 9 ...$

 For $f(6)$: $7: \boxed{1},\hspace{0.085in} 77: 2,\hspace{0.085in} 777: 3,\hspace{0.085in} 7777: 4,\hspace{0.085in} 77777: 6,\hspace{0.085in} 777777: 9 ...$

 Remember that it loops every $r$ times; the boxed numbers are where $N$ falls. Thus, the desired sum is $2+1+1+3+1 = \boxed{\textbf{(A) }8}$




\textbf{Solution 7 (Logs)}

Let's express $N$ as a geometric series:
 \[7+70+700+...+7 \cdot 10^{313}=N\]We can now apply the geometric series formula, thus:
 \[N=7 \cdot \frac{10^{313}-1}{9}\]Since we are dealing with large values, we can discard $-1$, as the value will virtually remain the same. Thus, 
 \[N=7 \cdot \frac{10^{313}}{9}\]Next, we must take the base $10$ logarithm of $N$, as we are trying to express the $rth$ roots of $N$ in the form $a \cdot 10^b$, thus 
 \[\log{\frac{7}{9} \cdot 10^{313}}=\]
 \[313 + \log{\frac{7}{9}}\]
 \[\approx 312.9\]So now to find the leading digits of the $rth$ roots of $N$, simply take $\log{\sqrt[r]{N}}=\frac{\log{N}}{r}=\frac{312.9}{r}$. This calculates the power of $10$ when $f(r)$ is taken in scientific notation ($a \cdot 10^b$). $a$ here in the notation corresponds with the fractional part of $\frac{312.9}{r}$ raised to the $10th$ power. We must find $a$ as it indicates the leading digit. Thus, for $f(2)$, 
 \[\frac{312.9}{2} = 156.45\]
 \[a=10^{0.45}\]Since $2<a<3$, $f(2)=2$ 
For $f(3)$, 
 \[\frac{312.9}{3} = 104.3\]
 \[a=10^{0.3}\]Since $1<a<2$, $f(3)=1$ 
For $f(4)$,
 \[\frac{312.9}{4} = 78.225\]
 \[a=10^{0.225}\]Since $1<a<2$, $f(4)=1$ 
For $f(5)$, 
 \[\frac{312.9}{5}=62.58\]
 \[a=10^{0.58}\]Since $3<a<4$, $f(5)=3$ 
For $f(6)$,
 \[\frac{312.9}{6}=52.15\]
 \[a=10^{0.15}\],Thus, $f(6)=1$So, $2+1+1+3+1=\boxed{8}$ 
(Note that this is very bashy, would not recommend on real test, as logs basically require a calculator)

 


 


	\item (\textit{AMC 10 10B Fall 2021}) In a particular game, each of $4$ players rolls a standard $6{ }$-sided die. The winner is the player who rolls the highest number. If there is a tie for the highest roll, those involved in the tie will roll again and this process will continue until one player wins. Hugo is one of the players in this game. What is the probability that Hugo's first roll was a $5,$ given that he won the game?


 $(\textbf{A})\: \frac{61}{216}\qquad(\textbf{B}) \: \frac{367}{1296}\qquad(\textbf{C}) \: \frac{41}{144}\qquad(\textbf{D}) \: \frac{185}{648}\qquad(\textbf{E}) \: \frac{11}{36}$

 \textbf{Answer: }C



\textbf{Solution 1}

Define the events $H$, $A$, and $B$ as follows:

 $H$ - Hugo wins.

 $A$ - Hugo's first roll is a $5$.

 $B$ - The highest number any player rolls on their first turn is $5$.

 The desired probability is $P(A|H)$.  In all games Hugo wins, the number Hugo rolls on his first turn must equal the highest number any player rolls on their first turn, so $P(A|H) = P(B|H)$.  By symmetry, $B$ and $H$ are independent events and $P(B|H) = P(B)$.  

 Define the random variable $X$ as the highest number any player rolls on their first turn.  

 Then $P(A|H) = P(B) = P(X = 5) = P(X \leq 5) - P(X \leq 4) = \left(\frac56\right)^4-\left(\frac46\right)^4=\frac{5^4-4^4}{6^4}=\frac{369}{1296}=\boxed{(\textbf{C}) \: \frac{41}{144}}.$




\textbf{Solution 2 (Conditional Probability)}

The conditional probability formula states that $P(A|B) = \frac{P(A\cap B)}{P(B)}$, where $A|B$ means A given B and $A\cap B$ means A and B. Therefore the probability that Hugo rolls a five given he won is $\frac{P(A \cap B)}{P(B)}$, where A is the probability that he rolls a five and B is the probability that he wins. In written form, 
 \[\text{P(Hugo rolled a 5 given he won)}=\frac{\text{P(Hugo rolls a 5 and wins)}}{\text{P(Hugo wins)}}.\]

 The probability that Hugo wins is $\frac{1}{4}$ by symmetry since there are four people playing and there is no bias for any one player. The probability that he gets a 5 and wins is more difficult; we will have to consider cases on how many players tie with Hugo. We consider the cases in which Hugo wins first, find the probabilities of those, and then account for the probability that Hugo rolled a 5, at the end.

 
$\textbf{Case 1:}$ No Players Tie With Hugo

 In this case, all other players must have numbers from 1 through four, and Hugo had a probability of 1 to win after the first round.

 There is a $\left(\frac{4}{6}\right)^{3} \cdot 1 = \frac{8}{27}$ chance of this happening.

 
$\textbf{Case 2:}$ One Player Ties With Hugo

 In this case, there are $3 \choose 1$ $= 3$ ways to choose which other player ties with Hugo, and the probability that this happens is $\frac{1}{6} \cdot \left(\frac{4}{6}\right)^2$. The probability that Hugo wins the game is then $\frac{1}{2}$ because there are now two players rolling the die. 

 Therefore the total probability in this case is $3 \cdot \frac{1}{2} \cdot \frac{1}{6} \cdot \left(\frac{4}{6}\right)^2= \frac{1}{9}$.

 $\textbf{Case 3:}$ Two Players Tie With Hugo

 In this case, there are $3 \choose 2$ $= 3$ ways to choose which other players tie with Hugo, and the probability that this happens is $\left(\frac{1}{6}\right)^2 \cdot \frac{4}{6}$. The probability that Hugo wins the game is then $\frac{1}{3}$ because there are now three players rolling the die. 

 Therefore the total probability in this case is $3 \cdot \frac{1}{3} \cdot \left(\frac{1}{6}\right)^2 \cdot \frac{4}{6} = \frac{1}{54}$.

 $\textbf{Case 4:}$ All Three Players Tie With Hugo

 In this case, the probability that all three players tie with Hugo is $\left(\frac{1}{6}\right)^3$. The probability that Hugo wins the game is $\frac{1}{4}$, so the total probability is $\frac{1}{4} \cdot \left(\frac{1}{6}\right)^3 = \frac{1}{864}$.

 Hugo has a $\frac{1}{6}$ probability of rolling a five himself, so the total probability that Hugo rolls a 5 and wins is

 
 \[\frac{1}{6}\left(\frac{8}{27} + \frac{1}{9} + \frac{1}{54} + \frac{1}{864}\right) = \frac{1}{6}\left(\frac{369}{864}\right) = \frac{1}{6}\left(\frac{41}{96}\right).\]

 Finally, the total probability is this probability divided by $\frac{1}{4}$ which is this probability times four; the final answer is

 
 \[4 \cdot \frac{1}{6}\left(\frac{41}{96}\right) = \frac{2}{3} \cdot \frac{41}{96} = \frac{41}{48 \cdot 3} = \frac{41}{144} = \boxed{C}.\]




\textbf{Solution 3}

We use $H$ to refer to Hugo.We use $H_1$ to denote the outcome of Hugo's $t$th toss.We denote by $A$, $B$, $C$ the other three players.We denote by $N$ the number of players among $A$, $B$, $C$ whose first tosses are 5.We use $W$ to denote the winner.

 We have
 \begin{align*} P \left( H_1 = 5 | W = H \right) & = \frac{P \left( H_1 = 5 , W = H \right)}{P \left( W = H \right)} \\ & = \frac{P \left( H_1 = 5 \right) P \left( W = H | H_1 = 5 \right) }{P \left( W = H \right)} \\ & = \frac{\frac{1}{6} P \left( W = H | H_1 = 5 \right)}{\frac{1}{4}} \\ & = \frac{2}{3} P \left( W = H | H_1 = 5 \right) . \end{align*}


 Now, we compute $P \left( W = H | H_1 = 5 \right)$.

 We have
 \begin{align*} & P \left( W = H | H_1 = 5 \right) \\ & = P \left( W = H | H_1 = 5 , \max \left\{ A_1, B_1, C_1 \right\} \leq 4 \right) P \left( \max \left\{ A_1, B_1, C_1 \right\} \leq 4 | H_1 = 5 \right) \\ & \quad + P \left( W = H | H_1 = 5 , \max \left\{ A_1, B_1, C_1 \right\} = 6 \right) P \left( \max \left\{ A_1, B_1, C_1 \right\} = 6 | H_1 = 5 \right) \\ & \quad + \sum_{N = 1}^3 P \left( W = H | H_1 = 5 , \max \left\{ A_1, B_1, C_1 \right\} = 5 , N \right) P \left( \max \left\{ A_1, B_1, C_1 \right\} = 5 , N  | H_1 = 5 \right) \\ & = P \left( W = H | H_1 = 5 , \max \left\{ A_1, B_1, C_1 \right\} \leq 4 \right) P \left( \max \left\{ A_1, B_1, C_1 \right\} \leq 4  \right) \\ & \quad + P \left( W = H | H_1 = 5 , \max \left\{ A_1, B_1, C_1 \right\} = 6 \right) P \left( \max \left\{ A_1, B_1, C_1 \right\} = 6 \right) \\ & \quad + \sum_{N = 1}^3 P \left( W = H | H_1 = 5 , \max \left\{ A_1, B_1, C_1 \right\} = 5 , N \right) P \left( \max \left\{ A_1, B_1, C_1 \right\} = 5 , N \right) \\ & = 1 \cdot P \left( \max \left\{ A_1, B_1, C_1 \right\} \leq 4  \right) + 0 \cdot P \left( \max \left\{ A_1, B_1, C_1 \right\} = 6 \right) \\ & \quad + \sum_{N = 1}^3 P \left( W = H | H_1 = 5 , \max \left\{ A_1, B_1, C_1 \right\} = 5 , N \right) P \left( \max \left\{ A_1, B_1, C_1 \right\} = 5 , N \right) \\ & = P \left( \max \left\{ A_1, B_1, C_1 \right\} \leq 4  \right) \\ & \quad + \sum_{N = 1}^3 P \left( W = H | H_1 = 5 , \max \left\{ A_1, B_1, C_1 \right\} = 5 , N \right) P \left( \max \left\{ A_1, B_1, C_1 \right\} = 5 , N \right) \\ & = \left( \frac{4}{6} \right)^3 + \sum_{N = 1}^3 \frac{1}{N + 1} \cdot \binom{3}{N} \left( \frac{1}{6} \right)^N \left( \frac{4}{6} \right)^{3 - N} \\ & = \frac{41}{96} . \end{align*}
The first equality follows from the law of total probability.The second equality follows from the property that Hugo's outcome is independent from other players' outcomes.

 Therefore,
 \begin{align*} P \left( H_1 = 5 | W = H \right) & = \frac{2}{3} P \left( W = H | H_1 = 5 \right) \\ & = \frac{2}{3} \frac{41}{96} \\ & = \frac{41}{144} , \end{align*}
so the answer is $\boxed{\textbf{(C) }\frac{41}{144}}$.

 


	\item (\textit{AMC 10 10B Fall 2021}) Regular polygons with $5,6,7,$ and $8$ sides are inscribed in the same circle. No two of the polygons share a vertex, and no three of their sides intersect at a common point. At how many points inside the circle do two of their sides intersect?


 $(\textbf{A})\: 52\qquad(\textbf{B}) \: 56\qquad(\textbf{C}) \: 60\qquad(\textbf{D}) \: 64\qquad(\textbf{E}) \: 68$

 \textbf{Answer: }E



\textbf{Solution}

Imagine we have $2$ regular polygons with $m$ and $n$ sides and $m>n$ inscribed in a circle without sharing a vertex. We see that each side of the polygon with $n$ sides (the polygon with fewer sides) will be intersected twice.(We can see this because to have a vertex of the $m$-gon on an arc subtended by a side of the $n$-gon, there will be one intersection to “enter” the arc and one to “exit” the arc. ~KingRavi)

 This means that we will end up with $2$ times the number of sides in the polygon with fewer sides.

 If we have polygons with $5,$ $6,$ $7,$ and $8$ sides, we need to consider each possible pair of polygons and count their intersections. 

 Throughout $6$ of these pairs, the $5$-sided polygon has the least number of sides $3$ times, the $6$-sided polygon has the least number of sides $2$ times, and the $7$-sided polygon has the least number of sides $1$ time.

 Therefore the number of intersections is $2\cdot(3\cdot5+2\cdot6+1\cdot7)=\boxed{(\textbf{E}) \:68}$.

 \textbf{Remark} For regular polygons with $5,6,7,$ and $8$ sides, the $68$ points of intersection inside the circle are shown below:
                    \begin{figure}[H]
                        \centering
                        \includegraphics[width=0.5\textwidth]{Solution/AMC_10_10B_Fall_2021_21_Solution_Figure_1.png}
                    \end{figure}
                    ~MRENTHUSIASM

 


	\item (\textit{AMC 10 10B Fall 2021}) For each integer $n\geq 2$, let $S_n$ be the sum of all products $jk$, where $j$ and $k$ are integers and $1\leq j<k\leq n$. What is the sum of the 10 least values of $n$ such that $S_n$ is divisible by $3$?


 $\textbf{(A)}\ 196 \qquad\textbf{(B)}\ 197 \qquad\textbf{(C)}\ 198 \qquad\textbf{(D)}\ 199 \qquad\textbf{(E)}\ 200$

 \textbf{Answer: }B



\textbf{Solution 1}

To get from $S_n$ to $S_{n+1}$, we add $1(n+1)+2(n+1)+\cdots +n(n+1)=(1+2+\cdots +n)(n+1)=\frac{n(n+1)^2}{2}$.

 Now, we can look at the different values of $n$ mod $3$. For $n\equiv 0\pmod{3}$ and $n\equiv 2\pmod{3}$, then we have $\frac{n(n+1)^2}{2}\equiv 0\pmod{3}$. However, for $n\equiv 1\pmod{3}$, we have
 \[\frac{1\cdot {2}^2}{2}\equiv 2\pmod{3}.\]

 Clearly, $S_2\equiv 2\pmod{3}.$ Using the above result, we have $S_5\equiv 1\pmod{3}$, and $S_8$, $S_9$, and $S_{10}$ are all divisible by $3$. After $3\cdot 3=9$, we have $S_{17}$, $S_{18}$, and $S_{19}$ all divisible by $3$, as well as $S_{26}, S_{27}, S_{28}$, and $S_{35}$. Thus, our answer is $8+9+10+17+18+19+26+27+28+35=27+54+81+35=162+35=\boxed{\mathrm{(B)}\ 197} .$ -BorealBear

 -minor edit by \href{https://artofproblemsolving.com/wiki/index.php?title=User:Yiyj1&action=edit&redlink=1}{Yiyj1}

 



\textbf{Solution 2 (bash)}

Since we have a wonky function, we start by trying out some small cases and see what happens. If $j$ is $1$ and $k$ is $2$, then there is one case. We have $2$ mod $3$ for this case. If $N$ is $3$, we have $1 \cdot 2 + 1 \cdot 3 + 2 \cdot 3$ which is still $2$ mod $3$. If $N$ is $4$, we have to add $1 \cdot 4 + 2 \cdot 4 + 3 \cdot 4$ which is a multiple of $3$, meaning that we are still at $2$ mod $3$. If we try a few more cases, we find that when $N$ is $8$, we get a multiple of $3$. When $N$ is $9$, we are adding $0$ mod $3$, and therefore, we are still at a multiple of $3$.

 When $N$ is $10$, then we get $0$ mod $3$ + $10(1+2+3+...+9)$ which is $10$ times a multiple of $3$. Therefore, we have another multiple of $3$. When $N$ is $11$, so we have $2$ mod $3$. So, every time we have $-1$ mod $9$, $0$ mod $9$, and $1$ mod $9$, we always have a multiple of $3$. Think about it: When $N$ is $1$, it will have to be $0 \cdot 1$, so it is a multiple of $3$. Therefore, our numbers are $8, 9, 10, 17, 18, 19, 26, 27, 28, 35$. Adding the numbers up, we get $\boxed{\textbf{(B) 197}}$




\textbf{Solution 3}

Denote $A_{n, <} = \left\{ \left( j , k \right) : 1 \leq j < k \leq n \right\}$, $A_{n, >} = \left\{ \left( j , k \right) : 1 \leq k < j \leq n \right\}$ and $A_{n, =} = \left\{ \left( j , k \right) : 1 \leq j = k \leq n \right\}$.

 Hence, $\sum_{\left( j , k \right) \in A_{n,<}} jk = \sum_{\left( j , k \right) \in A_{n,>}} jk = S_n$.

 Therefore,
 \begin{align*} S_n & = \frac{1}{2} \left( \sum_{\left( j , k \right) \in A_{n,<}} jk = \sum_{\left( j , k \right) \in A_{n,>}} jk \right) \\ & = \frac{1}{2} \left( \sum_{1 \leq j, k \leq n} jk - {\left( j , k \right) \in A_{n,=}} jk \right) \\ & = \frac{1}{2} \left( \sum_{j=1}^n \sum_{k=1}^n jk - \sum_{j=1}^n j^2 \right) \\ & = \frac{1}{2} \left( \frac{n^2 \left( n + 1 \right)^2}{4} - \frac{n \left( n + 1 \right) \left( 2 n + 1 \right) }{6} \right) \\ & = \frac{\left( n - 1 \right) n \left( n + 1 \right) \left( 3 n + 2 \right)}{24} . \end{align*}


 Hence, $S_n$ is divisible by 3 if and only if $\left( n - 1 \right) n \left( n + 1 \right) \left( 3 n + 2 \right)$ is divisible by $24 \cdot 3 = 8 \cdot 9$.

 First, $\left( n - 1 \right) n \left( n + 1 \right) \left( 3 n + 2 \right)$ is always divisible by 8. Otherwise, $S_n$ is not even an integer.

 Second, we find conditions for $n$, such that $\left( n - 1 \right) n \left( n + 1 \right) \left( 3 n + 2 \right)$ is divisible by 9.

 Because $3 n + 2$ is not divisible by 3, it cannot be divisible by 9.

 Hence, we need to find conditions for $n$, such that $\left( n - 1 \right) n \left( n + 1 \right)$ is divisible by 9.This holds of $n \equiv 0, \pm 1 \pmod{9}$.

 Therefore, the 10 least values of $n$ such that $\left( n - 1 \right) n \left( n + 1 \right)$ is divisible by 9 (equivalently, $S_n$ is divisible by 3) are 8, 9, 10, 17, 18, 19, 26, 27, 28, 35.Their sum is 197.

 Therefore, the answer is $\boxed{\textbf{(B) }197}$.

 


	\item (\textit{AMC 10 10B Fall 2021}) Each of the $5{ }$ sides and the $5{ }$ diagonals of a regular pentagon are randomly and independently colored red or blue with equal probability. What is the probability that there will be a triangle whose vertices are among the vertices of the pentagon such that all of its sides have the same color?


 $(\textbf{A})\: \frac23\qquad(\textbf{B}) \: \frac{105}{128}\qquad(\textbf{C}) \: \frac{125}{128}\qquad(\textbf{D}) \: \frac{253}{256}\qquad(\textbf{E}) \: 1$

 \textbf{Answer: }D



\textbf{Solution 1}

Instead of finding the probability of a same-colored triangle appearing, let us find the probability that one does not appear.After drawing the regular pentagon out, note the topmost vertex; it has 4 sides/diagonals emanating outward from it. We do casework on the color distribution of these sides/diagonals.

 $\textbf{Case 1}$: all 4 are colored one color. In that case, all of the remaining sides must be of the other color to not have a triangle where all three sides are of the same color. We can correspondingly fill out each color based on this constraint, but in this case you will always end up with a triangle where all three sides have the same color by inspection. 

 $\textbf{Case 2}$: 3 are one color and one is the other. Following the steps from the previous case, you can try filling out the colors, but will always arrive at a contradiction so this case does not work either.

 $\textbf{Case 3}$: 2 are one color and 2 are of the other color. Using the same logic as previously, we can color the pentagon 2 different ways by inspection to satisfy the requirements. There are ${4\choose2}$ ways to color the original sides/diagonals and 2 ways after that to color the remaining ones for a total of $6\cdot 2 = 12$ ways to color the pentagon so that no such triangle has the same color for all of its sides. 

 These are all the cases, and there are a total of $2^{10}$ ways to color the pentagon. Therefore the answer is $1-\frac{12}{1024} = 1-\frac{3}{256} = \frac{253}{256} = \boxed{D}$




\textbf{Solution 2 (Ramsey's Theorem)}

This problem is related to a special case of \href{https://en.wikipedia.org/wiki/Ramsey\%27s\_theorem}{Ramsey's Theorem}, \href{https://en.wikipedia.org/wiki/Ramsey\%27s\_theorem\#R(3,\_3)\_=\_6}{R(3, 3) = 6}. Suppose we color every edge of a $6$ vertex complete graph $(K_6)$ with $2$ colors, there must exist a $3$ vertex complete graph $(K_3)$ with all it's edges in the same color. That is, $K_6$ with edges in $2$ colors contains a monochromatic $K_3$. For $K_5$ with edges in $2$ colors, a monochromatic $K_3$ does not always exist.

 This is a problem about the probability of a monochromatic $K_3$ exist in a $5$ vertex complete graph $K_5$ with edges in $2$ colors.

 Choose a vertex, it has $4$ edges. 

 $\textbf{Case 1}$: When $3$ or more edges are the same color, there must exist a monochromatic $K_3$. Suppose the color is red, as shown below. 

 
                    \begin{figure}[H]
                        \centering
                        \includegraphics[width=0.5\textwidth]{Solution/AMC_10_10B_Fall_2021_23_Solution_Figure_1.jpg}
                    \end{figure}
                    

 There is only $1$ way to color all the edges in the same color. There is $\binom{4}{3} = 4$ ways to color $3$ edges in the same color. There are $2$ colors. The probability of $3$ or more edges the same color is $\frac{(1 + 4) \cdot 2}{2^4} = \frac{5}{8}$. So the probability of containing a monochromatic $K_3$ is $\frac{5}{8}$.

 $\textbf{Case 2}$: When $2$ edges are the same color, graphs that does not contain a monochromatic $K_3$ can exist. The following diagram shows steps to obtain graphs that does not contain a monochromatic $K_3$.

 
                    \begin{figure}[H]
                        \centering
                        \includegraphics[width=0.5\textwidth]{Solution/AMC_10_10B_Fall_2021_23_Solution_Figure_2.jpg}
                    \end{figure}
                    

 There are $\binom{4}{2} = 6$ ways to choose $2$ edges with the same color. For the other $4$ vertices there are $\binom{4}{2} = 6$ edges among them, there are $2^6 = 64$ ways to color the edges. There are only $2$ cases without a monochromatic $K_3$. 

 So the probability without monochromatic $K_3$ is $\frac{2}{64} = \frac{1}{32}$.

 The probability with monochromatic $K_3$ is $1 - \frac{1}{32} = \frac{31}{32}$.

 From case 1 and case 2, the probability with monochromatic $K_3$ is $\frac{5}{8} + \left( 1 - \frac{5}{8} \right) \cdot \frac{31}{32} = \boxed{(\textbf{D}) \frac{253}{256}}$




\textbf{Solution 3 (Elementary Casework)}

We can break this problem down into cases based on the distribution of the two colors throughout the $5$ sides(not diagonals) of the pentagon. This idea comes from the fact that if you play around with the sides for a while and their colors, you will see that the interior diagonals actually depend largely on the colors of the five exterior sides. The total number of combinations is $2^{10} = 1024$. We will be counting the number of arrangements in which there are no triangles with all $3$ sides the same color, and we can subtract from the total(complementary counting). 

 
            \begin{itemize}
        \item Exterior sides are the five sides of the regular pentagon
        \item Interior Diagonals are the five diagonals inside the pentagon made by the connection of two non-adjacent vertices
            \end{itemize}
            Case $1$: The five exterior sides are all blue or all red.

 There are $2$ options for the color given in the title of this case, and for each color, only one option is given for the exterior(All $5$ sides the same color; either Red or Blue). Once you color the five sides a single color, it is simple to notice that all the $5$ interior diagonals must be in the other color which the sides are not. To make it more clear, if all of the $5$ exterior sides are blue, all of the $5$ interior diagonals must be red, and if all of the $5$ exterior sides are red, all of the $5$ interior diagonals must be blue. This gives us a total of $2$ (Choices of Colors) $\cdot 1$ (Configuration for particular color existing on all 5 exterior edges) $= 2$ ways. 

 Case $2$: Four of the five exterior sides are in one color, while the remaining one exterior side is in the other color.

 There are $2$ ways to choose which color occupies four exterior sides and which color occupies the remaining one. Either we can have four red and one blue, or four blue and one red, which are the $2$ ways mentioned above. When we calculated the total number of combinations, we took into account that each segment could either be red or blue which gives $2$ choices for each of $10$ segments yielding $2^{10}$. What you must understand is that when we calculated this total, we did not account for any rotations or reflections or any other symmetry. For the same reason, we must not account for any symmetry such as rotations or reflections when calculating the number of arrangements where there is not a triangle consisting of three sides of the same color. There are $5$ ways to arrange the one blue and four reds or one red and four blues on the $5$ exterior sides(draw it out to see), and once you start playing with our condition, that no triangle may have all three sides in the same color, you will see that this case actually yields zero solutions as there comes a triangle in the middle consisting of all three sides of the same color. Hence this case yields $0$ ways.

 Case $3$: Three of the five exterior sides in one color and the remaining two in another color. 

 There are two sub-cases, one in which the two exterior sides which are colored differently from the other three are adjacent, and the other case in which they are separated by one other exterior side.

 Subcase $1$: The case in which the two exterior sides colored differently from the other three are adjacent. 

 If you draw the five exterior sides with two colored in one color and the other three colored in a different color, you will see, that there are absolutely no ways to color the interior diagonals such that no triangle with all three sides the same color exists. This subcase yields $0$ ways.

 Subcase $2$: The case in which the two exterior sides colored differently from the other three are separated by another exterior side. There are $5$ arrangements for the exterior (draw and see for yourself) and the colors can be swapped with each other and since each of the exterior configurations force a particular interior configuration, this yields $2 \cdot 5=10$ ways. 

 In total, there are $10+2=12$ ways, such that no triangle has all three sides in the same color. This yields $1024-12=1012$ ways such that there is a triangle such that it has all three sides in the same color. Thus, the probability is: $\frac{1012}{2^{10}} = \boxed{(\textbf{D}) \frac{253}{256}}$~Shiva Kannan

 


	\item (\textit{AMC 10 10B Fall 2021}) A cube is constructed from $4$ white unit cubes and $4$ blue unit cubes. How many different ways are there to construct the $2 \times 2 \times 2$ cube using these smaller cubes? (Two constructions are considered the same if one can be rotated to match the other.)


 $(\textbf{A})\: 7\qquad(\textbf{B}) \: 8\qquad(\textbf{C}) \: 9\qquad(\textbf{D}) \: 10\qquad(\textbf{E}) \: 11$

 \textbf{Answer: }A



\textbf{Solution 1 (Graph Theory)}

This problem is about the relationships between the white unit cubes and the blue unit cubes, which can be solved by \href{https://en.wikipedia.org/wiki/Graph\_theory}{Graph Theory}. We use a \href{https://en.wikipedia.org/wiki/Planar\_graph}{Planar Graph} to represent the larger cube. Each vertex of the planar graph represents a unit cube. Each edge of the planar graph represents a shared face between $2$ neighboring unit cubes. Each face of the planar graph represents a face of the larger cube.

 Now the problem becomes a \href{https://en.wikipedia.org/wiki/Graph\_coloring}{Graph Coloring} problem of how many ways to assign $4$ vertices blue and $4$ vertices white with \href{https://en.wikipedia.org/wiki/Homeomorphism}{Topological Equivalence}. For example, in Figure $(1)$, as long as the $4$ blue vertices belong to the same planar graph face, the different planar graphs are considered to be topological equivalent by rotating the larger cube.

 
                    \begin{figure}[H]
                        \centering
                        \includegraphics[width=0.5\textwidth]{Solution/AMC_10_10B_Fall_2021_24_Solution_Figure_1.jpg}
                    \end{figure}
                    

 Here is how the $4$ blue unit cubes are arranged: 

 In Figure $(1)$: $4$ blue unit cubes are on the same layer (horizontal or vertical).

 In Figure $(2)$: $4$ blue unit cubes are in $T$ shape.

 In Figure $(3)$ and $(4)$: $4$ blue unit cubes are in $S$ shape.

 In Figure $(5)$: $3$ blue unit cubes are in $L$ shape, and the other is isolated without a shared face.

 In Figure $(6)$: $2$ pairs of neighboring blue unit cubes are isolated from each other without a shared face.

 In Figure $(7)$: $4$ blue unit cubes are isolated from each other without a shared face.

 So the answer is $\boxed{\textbf{(A)}\ 7}$




\textbf{Solution 2 (Casework, Counting Down)}

Let’s split the cube into two layers; a bottom and top. Note that there must be four of each color, so however many number of one color are in the bottom, there will be four minus that number of the color on the top. We do casework on the color distribution of the bottom layer.

 \textbf{Case 1:} 4, 0

 In this case, there is only one possibility for the top layer - all of the other color - $\binom{4}{4}$. Therefore there is 1 construction from this case.

 \textbf{Case 2:} 3, 1

 In this case, the top layer has four possibilities, because there are four different ways to arrange it so that it also has a 3, 1 color distribution - $\binom{4}{3}$. Therefore there are 4 constructions from this case.

 \textbf{Case 3:} 2, 2

 In this case, the top layer has six possibilities of arrangement - $\binom{4}{2}$. However, having adjacent colors one way can be rotated to having adjacent colors any other way, so there is only one construction for the adjacent colors subcase and similarly, only one for the diagonal color subcase. Therefore the total number of constructions for this case is 2.

 The total number of constructions for the cube is thus $1+4+2=7=\boxed{A}$




\textbf{Solution 3 (Casework, Counting Up)}

Divide the $2 \times 2 \times 2$ cube into two layers, say, front and back. Any possible construction can be rotated such that the front layer has the same or greater number of white cubes than blue cubes, so we only need to count the number of cases given that is true.  

 
            \begin{enumerate}
        \item \textbf{Case 1: Each layer contains 2 cubes of each color.} Note that we only need to consider the configuration of the white cubes because all the other cubes will be blue cubes. There are 2 ways that the 2 white cubes in each layer can be arranged: adjacent or diagonal to each other.<ol><li>\textbf{Case 1.1: Both layers have 2 white cubes adjacent to each other.} Rotate the cube such that there are white cubes along the top edge of the front layer. Now, the white cubes in the back layer can be along the top, bottom, right, or left edges. So, case 1.1 results in $4$ constructions.</li><li>\textbf{Case 1.2: One layer has 2 white cubes adjacent to each other, and the other has 2 white cubes diagonal from each other.} Rotate the cube such that there are white cubes along the top of the front layer. The white cubes in the back layer can be at the top-left and bottom-right or at the top-right and bottom-left. If we rotate the former case by 90 degrees clockwise into the page, it becomes the same as the latter case. So, case 1.2 results in $1$ additional construction.</li><li>\textbf{Case 1.3: Both layers have white cubes diagonal from each other.} Rotate the cube such that there is a white cube at the top-left and bottom-right of the front layer. The back layer could also have white cubes at the top-left and bottom-right, but this is the same as case 1.1 with the white cubes in the back layer along the bottom edge. Alternatively, the back layer could have white cubes at the top-right and the bottom-left. This is a distinct case. So, case 1.3 results in $1$ additional construction.</li><li>So, case 1 results in $4+1+1=6$ distinct constructions.</li></ol>
        \item \textbf{Case 2: The front layer contains 3 white cubes, and the back layer contains one white cube.} In this case, unless the sole white and sole blue cubes in the front and back layers are on opposite corners of the $2 \times 2 \times 2$ cube, then that cube can be split into left and right layers with 2 cubes of each color in each layer. Those constructions were counted in case 1, so case 2 results in $1$ additional construction.
        \item \textbf{Case 3: The front layer contains 4 white cubes, and the back layer contains no white cubes.} If we split this construction into its left and right layers, then each layer will have 2 cubes of each color. So, this construction was counted in case 1, and case 3 results in $0$ additional constructions.
            \end{enumerate}
            Therefore, our answer is $6+1+0=\boxed{\textbf{(A)}\ 7}$.

 



\textbf{Solution 4 (Burnside Lemma)}

Burnside lemma is used to counting number of orbit where the element on the same orbit can be achieved by the defined operator, naming rotation, reflection and etc. 

 The facts for Burnside lemma are 

 1. the sum of stabilizer on the same orbit equals to the \# of operators; 

 2. the sum of stabilizer can be counted as $fix(g)$ 

 3. the sum of the $fix(g)/|G|$ equals the \# of orbit. 

 
Let's start with defining the operator for a cube, 

 1. $\textbf{e (identity)}$

 For identity, there are $\frac{8!}{4!4!} = 70$ 

 
2. ${\textbf r^{1}, r^{2}, r^{3}}$ to be the rotation axis along three pair of opposite face, 

 each contains $r^{i}_{90}, r^{i}_{180}, r^{i}_{270}$ where $i= 1, 2, 3$
$fix(r^{i}_{90}) = fix(r^{i}_{270}) = 2\cdot1 = 2$
$fix(r^{i}_{180}) = \frac{4!}{2!\cdot2!} = 6$

 therefore $fix(\textbf r^{i}) = 2+2+6 = 10$, and $fix(\textbf r^{1})+fix(\textbf r^{2})+fix(\textbf r^{3}) = 30$

 
3. ${\textbf r^{4}, r^{5}, r^{6}, r^{7}}$ to the rotation axis along four cube diagonals. 

 each contains $r^{i}_{120}, r^{i}_{240}$   where $i= 4, 5, 6, 7$
$fix(r^{i}_{120}) = fix(r^{i}_{240}) = 2\cdot1\cdot2\cdot1 = 4$

 therefore $fix(\textbf r^{i}) = 4+4 = 8$, and $fix(\textbf r^{4})+fix(\textbf r^{5})+fix(\textbf r^{6})+fix(\textbf r^{7}) = 32$

 


 4. ${\textbf r^{8}, r^{9}, r^{10}, r^{11}, r^{12}, r^{13}}$ to be the rotation axis along 6 pairs of diagonally opposite sides

 each contains $r^{i}_{180}$   where $i= 8, 9, 10, 11, 12, 13$
$fix(r^{i}_{180}) = \frac{4!}{2!\cdot2!} = 6$

 therefore $fix(\textbf r^{8})+fix(\textbf r^{9})+fix(\textbf r^{10})+fix(\textbf r^{11})+fix(\textbf r^{12})+fix(\textbf r^{13}) = 36$

 
5. The total number of operators are 

 $|G| = 1 + 3\cdot3 + 4\cdot2 + 6\cdot1 = 24$

 Based on 1, 2, 3, 4 the total number of stabilizers is $70 + 30 + 32 + 36 = 168$

 therefore the number of orbits $= \frac{168}{G=24} = \boxed{7}$




\textbf{Solution 5 (Quick 'n Easy but not rigorous)}

Since rotations of a single pattern are considered indistinguishable, we can assume that the forward upper right corner of the 2-by-2-by-2 cube is a blue cube (since we can always rotate the big cube to place a blue cube in that spot).

 Once we've assigned this cube to be blue, we note that 3 1-by-1-by-1 cubes share a side with it, 3 1-by-1-by-1 cubes share a corner with it, and 1 1-by-1-by-1 cube does not touch the assigned cube at all, from the perspective of someone who can only see the cube's faces. We'll call the first 3 "adjacent", the second 3 "cornering", and the last one "opposite."

 We can use a little bit of intuition to confirm that due to the rotation condition, we should treat all adjacents as indistinguishable, all cornerings as indistinguishable, and of course the opposite one is unique from all the others. Thus, we can list out like so (keeping in mind that there are 3 adjacents, 3 cornerings, and 1 opposite, and that we're choosing the positions of the remaining 3 blue cubes):

 OAA, OAC, OCC, CCC, CAA, CCA, AAA.

 This gives the answer to be $\boxed{\textbf{(A)}\ 7}$.

 


	\item (\textit{AMC 10 10B Fall 2021}) A rectangle with side lengths $1{ }$ and $3,$ a square with side length $1,$ and a rectangle $R$ are inscribed inside a larger square as shown. The sum of all possible values for the area of $R$ can be written in the form $\frac mn$, where $m$ and $n$ are relatively prime positive integers. What is $m+n?$

                    \begin{figure}[H]
                        \centering
                        \includegraphics[width=0.5\textwidth]{Problem/AMC_10_10B_Fall_2021_25_Problem_Figure_1.png}
                    \end{figure}
                    
$(\textbf{A})\: 14\qquad(\textbf{B}) \: 23\qquad(\textbf{C}) \: 46\qquad(\textbf{D}) \: 59\qquad(\textbf{E}) \: 67$

 \textbf{Answer: }E



\textbf{Solution 1}

We see that the polygon bounded by the small square, large square, and rectangle of known lengths is an isosceles triangle. Let’s draw a perpendicular from the vertex of this triangle to its opposing side;

 
                    \begin{figure}[H]
                        \centering
                        \includegraphics[width=0.5\textwidth]{Solution/AMC_10_10B_Fall_2021_25_Solution_Figure_1.png}
                    \end{figure}
                    

 We see that this creates two congruent triangles. Let the smaller side of the triangle have length $a$ and let the larger side of the triangle have length $b$. Now we see by AAS congruency that if we draw perpendiculars that surround the smaller square, each outer triangle will be congruent to these two triangles.

 Now notice that these small triangles are also similar to the large triangle bounded by the bigger square and the rectangle by AA, and the ratio of the sides are 1:3, so we can fill in the lengths of that triangle. Similarly, the small triangle on the right bounded by the rectangle and the square is also congruent to the other small triangles by AAS, so we can fill in those sides;

 
                    \begin{figure}[H]
                        \centering
                        \includegraphics[width=0.5\textwidth]{Solution/AMC_10_10B_Fall_2021_25_Solution_Figure_2.png}
                    \end{figure}
                    

 Since the larger square by definition has all equal sides, we can set the sum of the lengths of the sides equal to each other.$3a+b+b+a = 3b+a \implies 3a = b$. Now let's draw some more perpendiculars and rename the side lengths.

 
                    \begin{figure}[H]
                        \centering
                        \includegraphics[width=0.5\textwidth]{Solution/AMC_10_10B_Fall_2021_25_Solution_Figure_3.png}
                    \end{figure}
                    

 By AA similarity, when we draw a perpendicular from the intersection of the two rectangles to the large square, we create a triangle below that is similar to the small congruent triangles with length $a,3a$. Since we don't know its scale, we'll label its sides $c,3c$.

 The triangle that is created above the perpendicular is congruent to the triangle on the opposite of the rectangle with unknown dimensions because they share the same hypotenuse and have two angles in common. Thus we can label these two triangles accordingly.

 The side length of the big square is $10a$, so we can find the remaining dimensions of the triangle bounded by the rectangle with unknown dimensions and the large square in terms of $a$ and $c$:

 
                    \begin{figure}[H]
                        \centering
                        \includegraphics[width=0.5\textwidth]{Solution/AMC_10_10B_Fall_2021_25_Solution_Figure_4.png}
                    \end{figure}
                    

 This triangle with side lengths $4a-c$ and $6a-3c$ is similar to the triangle directly below it with side lengths $3a$ and $3c$ by AA similarity, so we can set up a ratio equation:$\frac{3a}{3c} = \frac{6a-3c}{4a-c} \implies 4a^2-ac = -3c^2 + 6ac \implies 4a^2 - 7ac + 3c^2 = 0 \implies (4a-3c)(a-c) = 0$.There are two solutions to this equation; $c = \frac{4}{3}a$ and $c = a$. For the first solution, the triangle in the corner has sides $2a$ and $\frac{8}{3}a$. Using Pythagorean theorem on that triangle, the hypotenuse has length $\frac{10}{3}a$.The triangle directly below has side lengths $3a$ and $4a$ in this case, so special right triangle yields the hypotenuse to be $5a$.The area of the rectangle is thus $5a\cdot\frac{10}{3}a = \frac{50}{3}a^2$.For the second solution, the side lengths of the corner triangle are $3a$ and $3a$, so the hypotenuse of the triangle is $3\sqrt{2}a$. The triangle below that also has side lengths $3a$ and $3a$, so its hypotenuse is the same. Then the area of the rectangle is $(3\sqrt{2}a)^2 = 18a^2$.

 The sum of the possible areas of the rectangle is therefore $18a^2+\frac{50}{3}a^2 = \frac{104}{3}a^2$.

 Using Pythagorean theorem on the original small congruent triangles, $a^2+9a^2 = 1$ or $a^2 = \frac{1}{10}$. Therefore the sum of the possible areas of the rectangle is $\frac{104}{3}\cdot\frac{1}{10} = \frac{52}{15}$. Therefore $m = 52$, $n = 15$, and $m + n = 67 = \boxed{E}$




\textbf{Solution 2}

We use \href{https://artofproblemsolving.com/wiki/index.php?title=File:2021\_AMC\_10B\_(Nov)\_Problem\_25,\_sol.png}{Image:2021\_AMC\_10B\_(Nov)\_Problem\_25,\_sol.png} to facilitate our analysis.

 Denote $\angle AFE = \theta$. Thus, $\angle FIB = \angle CEF = \angle EKG = \angle KLC = \theta$.

 Hence,
 \begin{align*} AB & = AF + FB \\ & = EF \cos \angle EFA + IF \sin \angle FIB \\ & = 3 \cos \theta + \sin \theta , \end{align*}
and
 \begin{align*} AC & = AE + EK + KC \\ & = EF \sin \angle EFA + EG \cos \angle CEG + KG \cos \angle EKG + KL \sin \angle CLK \\ & = 3 \sin \theta + \cos \theta + \cos \theta + \sin \theta \\ & = 2 \cos \theta + 4 \sin \theta . \end{align*}


 Because $AB = AC$, $3 \cos \theta + \sin \theta = 2 \cos \theta + 4 \sin \theta$.Hence, $\tan \theta = \frac{1}{3}$.Hence, $\sin \theta = \frac{1}{\sqrt{10}}$ and $\cos \theta = \frac{3}{\sqrt{10}}$.

 Hence, $AB = AC = BD = CD = \sqrt{10}$.

 Now, we put the graph to a coordinate plane by setting point $A$ as the origin, putting $AB$ in the $x$-axis and $AC$ on the $y$-axis.

 Hence, $A = \left( 0 , 0 \right)$, $B = \left( \sqrt{10} , 0 \right)$, $C = \left( 0 , \sqrt{10} \right)$, $D = \left( \sqrt{10} , \sqrt{10} \right)$, $E = \left( 0 , \frac{3}{\sqrt{10}} \right)$, $F = \left( \frac{9}{\sqrt{10}} , 0 \right)$, $G = \left( \frac{1}{\sqrt{10}} , \frac{6}{\sqrt{10}} \right)$, $H = \left( \frac{4}{\sqrt{10}} , \frac{7}{\sqrt{10}} \right)$, $I = \left( \sqrt{10} , \frac{3}{\sqrt{10}} \right)$.

 Denote $P = \left( \frac{10 - u}{\sqrt{10}} , \sqrt{10} \right)$, $Q = \left( \sqrt{10} , \frac{10 - v}{\sqrt{10}} \right)$.

 Because $HPQJ$ is a rectangle, $HP \perp PQ$. Hence, $m_{HP} m_{PQ} = -1$.We have $m_{HP} = \frac{3}{6 - u}$ and $m_{PQ} = - \frac{v}{u}$.Hence,
 \[ \frac{3}{6 - u} \cdot \left( - \frac{v}{u} \right) = - 1 . \hspace{1cm} (1) \]

 Because $HPQJ$ is a rectangle, $x_J + x_P = x_H + x_Q$ and $y_J + y_P = y_H + y_Q$.Hence, $J = \left( \frac{4 + u}{\sqrt{10}} , \frac{7 - v}{\sqrt{10}} \right)$.

 The equation of line $GI$ is
 \begin{align*} y & = \frac{\frac{3}{\sqrt{10}} - \frac{6}{\sqrt{10}}}{\sqrt{10} - \frac{1}{\sqrt{10}}} \left( x - \frac{1}{\sqrt{10}} \right) + \frac{6}{\sqrt{10}} \\ & = - \frac{x}{3} + \frac{19}{3 \sqrt{10}} . \end{align*}
Because point $J$ is on line $GI$, plugging the coordinates of $J$ into the equation of line $GI$, we get
 \[ \frac{7 - v}{\sqrt{10}} = - \frac{\frac{4 + u}{\sqrt{10}}}{3} + \frac{19}{3 \sqrt{10}} . \hspace{1cm} (2) \]

 By solving Equations (1) and (2), we get $\left( u , v \right) = \left( 2 , \frac{8}{3} \right)$ or $\left( 3 , 3 \right)$.

 Case 1: $\left( u , v \right) = \left( 2 , \frac{8}{3} \right)$.

 We have $P = \left( \frac{8}{\sqrt{10}} , \sqrt{10} \right)$ and $Q = \left( \sqrt{10} , \frac{22}{3 \sqrt{10}} \right)$.Thus, $HP = \frac{5}{\sqrt{10}}$ and $PQ = \frac{10}{3\sqrt{10}}$.

 Therefore, ${\textrm Area} \ R = HP \cdot PQ = \frac{5}{3}$.

 Case 2: $\left( u , v \right) = \left( 3 , 3 \right)$.

 We have $P = \left( \frac{7}{\sqrt{10}} , \sqrt{10} \right)$ and $Q = \left( \sqrt{10} , \frac{7}{ \sqrt{10}} \right)$.Thus, $HP = \frac{3 \sqrt{2}}{\sqrt{10}}$ and $PQ = \frac{3 \sqrt{2}}{\sqrt{10}}$.

 Therefore, ${\textrm Area} \ R = HP \cdot PQ = \frac{9}{5}$.

 Putting these two cases together, the sum of all possible values of the area of $R$ is $\frac{5}{3} + \frac{9}{5} = \frac{52}{15}$.

 
Therefore, the answer is $\boxed{\textbf{(E) }67}$.

 



\textbf{Solution 3}


                    \begin{figure}[H]
                        \centering
                        \includegraphics[width=0.5\textwidth]{Solution/AMC_10_10B_Fall_2021_25_Solution_Figure_5.png}
                    \end{figure}
                    We will scale every number up by a factor of $\sqrt{10}.$ This implies our final area will be $\frac{1}{10}$ of the answer we receive.

 We have
 \[FAE\sim EOH \sim IOH\sim JDI\sim KCJ \sim NQK\sim GPF.\]Let $AE=a$ and $FA=b.$ We have
 \[FP=AE=OH=JC=\frac13ID=a\]and
 \[PG=AF=EO=OI=KC=\frac13 DJ=b\]As $ABCD$ is a square, we have $AD=DC$ or
 \[a+2b+3a=3b+a \Rightarrow 3a=b.\]Since $a^2+b^2=10,$ we have
 \[a=1,b=3.\]We have $\triangle GPL\cong \triangle MQN$ which implies
 \[MQ=GP=3.\]Denote $QK=x.$ As $\triangle NQK\sim \triangle JDI,$ we have $NQ=3x.$

 We have
 \begin{align*}BM&=BC-(CK+QK+MQ)\\ &=4-x.\end{align*}


 In addition,
 \begin{align*}LB&=AB-(AF+FP+PL)\\&=6-3x.\end{align*}


 Since $\triangle LBM\sim \triangle MQN,$ we have
 \[\frac{LB}{BM}=\frac{MQ}{QN}\Rightarrow \frac{6-3x}{4-x}=\frac{3}{3x}=\frac{1}{x}.\]Simplifying we have
 \[3x^2-7x+4=0 \Rightarrow x=\frac43, 1.\]We have
 \begin{align*}[GLMN]&=MN\cdot LM\\ &= 3\sqrt{x^2+1}\cdot \sqrt{10x^2-44x+52}.\end{align*}


 Plugging in $x=1,$ we have $[GLMN]=18.$

 Plugging in $x=\frac43,$ we have $[GLMN]=\frac{50}{3}.$

 The sum of the two possible $R$s is
 \[\frac{1}{10}\cdot\frac{104}{3}=\frac{52}{15}.\]Hence, $52+15=\boxed{\textbf{(E) }67}.$




\textbf{Solution 4}

We use the same diagram as solution 3, without scaling by $\sqrt{10}$. Denote $S$ the foot from $G$ to $HN$. Using the same method, obtain $JC=\frac{1}{\sqrt{10}}$ and $KC=\frac{3}{\sqrt{10}}$. 

 We have $\angle GHS = 2 \angle HIO \Longrightarrow \sin \angle GHS = 2 \sin \angle HIO \cos \angle HIO = \frac{3}{5}$, so $GS = \frac{3}{5}$, and $HS=\frac{4}{5}$. 

 Denote $x=KN$, then $SN=HK-HS-NK=\frac{11-5x}{5}$, and $NQ=\frac{3x}{\sqrt{10}}$. 

 Also, $\triangle NQM \cong \triangle LPG$ because they're similar by $AA$, and then $LG=NM$, so $MQ=PG=\frac{3}{\sqrt{10}}$. 

 
Thus $\tan \angle KNQ = \frac{1}{3}$ and $\tan \angle QNM = \frac{1}{x}$. Then $\cot \angle GNS = \tan (90 ^{\circ} - \angle GNS) = \tan \angle KNM = \tan (\angle KNQ + \angle QNM) = \frac{\tan \angle KNQ + \tan \angle QNM}{1-\tan \angle KNQ \tan \angle QNM}$

 which simplifies to

 $\frac{11-5x}{3}=\frac{3+x}{3x-1}$

 Cross multiplying gives

 $3x^2-7x+4=0$, with solutions $x=1,\frac{4}{3}$. 

 Plugging these back into the setup gives areas of $\frac{9}{5}$ and $\frac{5}{3}$, respectively, which have sum $\frac{52}{15}$, and thus the answer is $52+15=\boxed{\textbf{(E) }67}.$



\end{enumerate}